\chapter*{导数及其运用}

\subsection*{导数的概念和意义(1)导数的概念}

A 组：

\begin{enumerate}
    \item 函数 $f\left( x\right)  = {x}^{3}$ 在区间 $\left\lbrack  {3,4}\right\rbrack$ 上的平均变化率等于\blkda{37}.
    \begin{jieda}
        由函数 $f\left( x\right)  = {x}^{3}$ ,可得函数 $f\left( x\right)$ 在 $\left\lbrack  {3,4}\right\rbrack$ 上的平均变化率为$\dfrac{\Delta y}{\Delta x} = \dfrac{f\left( 4\right)  - f\left( 3\right) }{4 - 3} = \dfrac{{4}^{3} - {3}^{3}}{4 - 3} = {37}$
    \end{jieda}
    
    \item 设函数 $f\left( x\right)  = {x}^{2} - 1$ ，当自变量 $x$ 由 1 变到 1.1 时，函数的平均变化率是\blkda{$2.1$}.
    \begin{jieda}
        $\dfrac{[(1.1)^{2} - 1]-[1^{2} - 1]}{1.1-1} = 2.1$
    \end{jieda}
    
    \item 设函数 $y = f\left( x\right)$ 在 $x = {x}_{0}$ 处可导,且 ${f}^{\prime }\left( {x}_{0}\right)  = 1$ ,则 $\mathop{\lim }\limits_{{h \rightarrow  0}}\dfrac{f\left( {{x}_{0} + {3h}}\right)  - f\left( {x}_{0}\right) }{h} =$ \blkda{3}.
    \begin{jieda}
        根据导数的定义可知: $\mathop{\lim }\limits_{{h \rightarrow  0}}\dfrac{f\left( {{x}_{0} + {3h}}\right)  - f\left( {x}_{0}\right) }{h} = 3\mathop{\lim }\limits_{{h \rightarrow  0}}\dfrac{f\left( {{x}_{0} + {3h}}\right)  - f\left( {x}_{0}\right) }{3h} = 3{f}^{\prime }\left( {x}_{0}\right)  = 3$.
    \end{jieda}
    
    \item 已知物体做直线运动的方程为 $s = s\left( t\right)$ ，则 ${s}^{\prime }\left( 4\right)  = {10}$ 表示的意义是\dotfill{(\kh{D})}
    \xx{经过 $4{s}$ 后物体向前走了 ${10}{m}$}{物体在前 4 秒内的平均速度为 ${10}{m}/{s}$}{物体在第 4 秒内向前走了 $10m$}{物体在第 4 秒末的瞬时速度为 ${10}{m}/{s}$}
    
    \item 函数 ${y} = {f}\left( {x}\right)$ 的图像如图所示. 求割线 PQ 的斜率； 当点 Q 沿曲线向 P 运动时，割线 PQ 的斜率会变大还是变小?
    
    \begin{center}
    \includegraphics[width=0.2\textwidth]{bodyxb/src/bo_d6fgfps601uc73e06260_0_271_896_290_245_0.jpg}
    \end{center}
    \begin{jieda}
        $k_{PQ} = \dfrac{4-1}{5-1} = \dfrac{3}{4}$；变大
    \end{jieda}

\end{enumerate}

B 组：
\begin{enumerate}
    \item 函数 $f\left( x\right)  = {x}^{2}$ 在区间 $\left\lbrack  {0,2}\right\rbrack$ 上的平均变化率等于 $x = m$ 时的瞬时变化率，则 $m =$\blkda{1}
    \begin{jieda}
        函数 $f\left( x\right)  = {x}^{2}$ 在区间 $\left\lbrack  {0,2}\right\rbrack$ 上的平均变化率为 $\dfrac{f\left( 2\right)  - f\left( 0\right) }{2 - 0} = \dfrac{4 - 0}{2 - 0} = 2$ ,$f\left( x\right)  = {x}^{2}$ 在 $x = m$ 时的瞬时变化率为
        $\mathop{\lim }\limits_{{{\Delta x} \rightarrow  0}}\dfrac{f\left( {m + {\Delta x}}\right)  - f\left( m\right) }{\Delta x} = \mathop{\lim }\limits_{{{\Delta x} \rightarrow  0}}\dfrac{{\left( m + \Delta x\right) }^{2} - {m}^{2}}{\Delta x} = \mathop{\lim }\limits_{{{\Delta x} \rightarrow  0}}\left( {{\Delta x} + {2m}}\right)  = {2m},$
        所以 $2 = {2m}$ ,解得 $m = 1$ .
    \end{jieda}
    
    \item 已知 $f\left( x\right)  =  - {2x} - 5$ 和 $g\left( x\right)  =  - {3x} + 7$ 在区间 $\left\lbrack  {m,n}\right\rbrack$ 上的平均变化率分别为 $a$ 和 $b$ ,则 $a + b =$ \blkda{$-5$}.
    \begin{jieda}
        由题意，$a = \dfrac{f\left( n\right)  - f\left( m\right) }{n - m} = \dfrac{-{2n} - 5 - \left( {-{2m} - 5}\right) }{n - m} =  - 2$
        
        $b = \dfrac{g\left( n\right)  - g\left( m\right) }{n - m} = \dfrac{-{3n} + 7 - \left( {-{3m} + 7}\right) }{n - m} =  - 3$，故 $a + b =  - 5$ .
    \end{jieda}
    
    \item 函数 $y = {x}^{2} + x$ 在 $x = 2$ 到 $x = 2 + h$ 之间的平均变化率为\blkda{$h+5$}.
    \begin{jieda}
        $\dfrac{[(2+h)^2+(2+h)] - (4+2)}{h} = h+5$
    \end{jieda}
    
    \item 若函数 $f\left( x\right)  = {x}^{2} - {m}^{2}$ 在区间 $\left\lbrack  {2,t}\right\rbrack$ 上的平均变化率为 5,则 $t =$ \blkda{3}.
    \begin{jieda}
        $\dfrac{(t^2-m^2) - (4-m^2)}{t-2} = 5$，解得$t = 3$
    \end{jieda}
    
    \item 已知函数 $f\left( x\right)  = \cos x$ ，则 $\mathop{\lim }\limits_{{h \rightarrow  0}}\dfrac{f\left( {\dfrac{\pi }{2} + h}\right)  - f\left( \dfrac{\pi }{2}\right) }{h} =$ \blkda{$-1$}.
    \begin{jieda}
        $\mathop{\lim }\limits_{{h \rightarrow  0}}\dfrac{f\left( {\dfrac{\pi }{2} + h}\right)  - f\left( \dfrac{\pi }{2}\right) }{h} = \mathop{\lim }\limits_{{h \rightarrow  0}}\dfrac{cos\left( {\dfrac{\pi }{2} + h}\right)  - cos\left( \dfrac{\pi }{2}\right) }{h} = \mathop{\lim }\limits_{{h \rightarrow  0}}\dfrac{-sin(h)}{h} = -1$
    \end{jieda}
    
    \item 已知 ${f}^{\prime }\left( {x}_{0}\right)  = m$ ，则 $\mathop{\lim }\limits_{{h \rightarrow  0}}\dfrac{f\left( {{x}_{0} - {3h}}\right)  - f\left( {x}_{0}\right) }{h} =$ \blkda{$-3m$}.
    \begin{jieda}
        根据导数的定义可知$\mathop{\lim }\limits_{{h \rightarrow  0}}\dfrac{f\left( {{x}_{0} - {3h}}\right)  - f\left( {x}_{0}\right) }{3h} = m$，故$\mathop{\lim }\limits_{{h \rightarrow  0}}\dfrac{f\left( {{x}_{0} - {3h}}\right)  - f\left( {x}_{0}\right) }{h} = -3m$.
    \end{jieda}
    
    \item 已知函数 $y = f\left( x\right)$ ，若 ${f}^{\prime }\left( 1\right)  = 2$ ，则 $\mathop{\lim }\limits_{{k \rightarrow  0}}\dfrac{f\left( {1 - k}\right)  - f\left( 1\right) }{2k} =$ \blkda{$-1$}.
    \begin{jieda}
        根据导数的定义可知$\mathop{\lim }\limits_{{k \rightarrow  0}}\dfrac{f\left( {1 - k}\right)  - f\left( 1\right) }{-k} = 2$，故$\mathop{\lim }\limits_{{k \rightarrow  0}}\dfrac{f\left( {1 - k}\right)  - f\left( 1\right) }{2k} = -1$.
    \end{jieda}
    
    \item 如图，在平面直角坐标系 ${xOy}$ 中，直线 $y = \sqrt{3}x,y = 0$ ， $x = t\left( {t > 0}\right)$ 围成的 $\bigtriangleup {OAB}$ 的面积为 $S\left( t\right)$ ,则 $S\left( t\right)$ 在 $t = 2$ 时的瞬时变化率是\blkda{$2\sqrt{3}$}.
    \begin{center}
    \includegraphics[width=0.2\textwidth]{bodyxb/src/bo_d6fgfps601uc73e06260_1_235_209_337_377_0.jpg}
    \end{center}
    \begin{jieda}
        由题设知: $B\left( {t,\sqrt{3}t}\right)$ ,即 ${AB} = \sqrt{3}t$ ,而 ${OA} = t$，所以 $S\left( t\right)  = \dfrac{1}{2}{OA} \cdot  {AB} = \dfrac{\sqrt{3}}{2}{t}^{2}$ ,则${S}^{\prime }\left( 2\right)  = \mathop{\lim }\limits_{{{\Delta t} \rightarrow  0}}\dfrac{S\left( {2 + {\Delta t}}\right)  - S\left( 2\right) }{\Delta t} = \mathop{\lim }\limits_{{{\Delta t} \rightarrow  0}}\dfrac{2\sqrt{3}{\Delta t} + \dfrac{\sqrt{3}}{2}\Delta {t}^{2}}{\Delta t} = \mathop{\lim }\limits_{{{\Delta t} \rightarrow  0}}\left( {2\sqrt{3} + \dfrac{\sqrt{3}}{2}{\Delta t}}\right)  = 2\sqrt{3}.$

    \end{jieda}
    
    \item 某跳水运动员在距离地面3m高的跳台上练习跳水，其重心相对于水面的高度 $h$ (单位: ${m}$ ) 与起跳后的时间 $t$ (单位: ${s}$ ) 的函数关系是 $h\left( t\right)  =  - {4.9}{t}^{2} + {2t} + 4$ , 则该运动员在 $t = {0.5}$ 时的瞬时速度(单位: ${m}/{s}$ )为\blkda{$-2.9$}.
    \begin{jieda}
        由题意 $h\left( t\right)  =  - {4.9}{t}^{2} + {2t} + 4$ ,求导后得 ${h}^{\prime }\left( t\right)  =  - {9.8t} + 2$ ,当 $t = {0.5}$ 时，${h}^{\prime }\left( {0.5}\right)  =  - {9.8} \times  {0.5} + 2 =  - {2.9}$
    \end{jieda}
    
    \item 已知在一次降雨过程中，某地降雨量 $y$ (单位: ${{mm}}$ )与时间(单位: ${{mm}}$ ) 的函数关系可近似表示为 $y = \sqrt{10t}$ ，则在 $t = {40}{{min}}$ 时的瞬时降雨强度(某一时刻降雨量的瞬间变化率)为\blkda{$\dfrac{1}{4}$} ${{mm}}/{{min}}$ .

    \begin{jieda}
        $\mathop{\lim }\limits_{{h \rightarrow  0}}\dfrac{\sqrt{10\times(40+h)}-\sqrt{10\times40}}{h} = \mathop{\lim }\limits_{{h \rightarrow  0}}\dfrac{10h}{h(\sqrt{10\times(40+h)}+\sqrt{10\times40})} = \mathop{\lim }\limits_{{h \rightarrow  0}}\dfrac{10}{(\sqrt{10\times(40+h)}+\sqrt{10\times40})} = \dfrac{1}{4}$
    \end{jieda}
    
    \item 如图，水缸为圆锥形，圆锥底面直径为 2.5 米，高为 5 米， 水被以 $\pi$ 立方米/小时的速度注入水缸中，当水缸中的水深为 2 米时，此时水面上升的瞬时速度(变化率)为\blkda{4}(单位:米/ 小时) .
    \begin{center}
    \includegraphics[width=0.2\textwidth]{bodyxb/src/bo_d6fgfps601uc73e06260_1_1199_1067_214_328_0.jpg}
    \end{center}
    
    \item 已知函数 $y = f\left( x\right)$ 的图象如图所示， ${f}^{\prime }\left( x\right)$ 是函数 $f\left( x\right)$ 的导函数，则\dotfill{(\kh{A})}

    \begin{center}
    \includegraphics[width=0.3\textwidth]{bodyxb/src/bo_d6fgfps601uc73e06260_1_237_1480_356_331_0.jpg}
    \end{center}

    \xx{${f}^{\prime }\left( 2\right)  < \dfrac{f\left( 4\right)  - f\left( 2\right) }{2} < {f}^{\prime }\left( 4\right)$}{${f}^{\prime }\left( 4\right)  < {f}^{\prime }\left( 2\right)  < \dfrac{f\left( 4\right)  - f\left( 2\right) }{2}$}{${f}^{\prime }\left( 2\right)  < {f}^{\prime }\left( 4\right)  < \dfrac{f\left( 4\right)  - f\left( 2\right) }{2}$}{$\dfrac{f\left( 4\right)  - f\left( 2\right) }{2} < {f}^{\prime }\left( 4\right)  < {f}^{\prime }\left( 2\right)$}
    
    \item 汽车行驶的路程 $S$ 和时间 $t$ 之间的函数图像如图所示,在时间段 $\left\lbrack  {{t}_{0},{t}_{1}}\right\rbrack$ , $\left\lbrack  {{t}_{1},{t}_{2}}\right\rbrack  ,\left\lbrack  {{t}_{2},{t}_{3}}\right\rbrack$ 上的平均速度分别为 $\overline{{v}_{1}},\overline{{v}_{2}},\overline{{v}_{3}}$ ,则三者的大小关系为\blkda[3]{$\overline{{v}_{1}}<\overline{{v}_{2}}<\overline{{v}_{3}}$}.
    
    \begin{center}
    \includegraphics[width=0.3\textwidth]{bodyxb/src/bo_d6fgfps601uc73e06260_2_645_436_354_293_0.jpg}
    \end{center}
    
    \item 自由落体运动中，物体下落的距离d(单位:m)与时间t(单位:s)近似满足函数关系 $d = 5{t}^{2}$ .
    \begin{enumerate}
        \item 求物体在$[2, 4]$时间段内的平均速度
        \item 求物体在 $t = 3$ 时的瞬时速度
        \item 求物体在 $t = a\left( {a > 0}\right)$ 时的瞬时速度
    \end{enumerate}
    \begin{jieda}
        \begin{enumerate}
            \item ${30}{m}/{s}$
            
            \item ${30}{m}/{s}$
            
            \item ${10a}{m}/{s}$
        \end{enumerate}
    \end{jieda}
    
    \item 自由落体运动的位移 ${d}$ (单位: ${m}$ ) 和时间 ${t}$ (单位: ${s}$ ) 满足函数关系 ${d} = \dfrac{1}{2}{g}{t}^{2}$ (g 为重力加速度).
    \begin{enumerate}
        \item 分别求 $\left\lbrack  {4,{4.1}}\right\rbrack  ,\left\lbrack  {4,{4.01}}\right\rbrack  ,\left\lbrack  {4,{4.001}}\right\rbrack$ 这些时间段内自由落体的平均速度
        \item 求 $t = 4$ 时的瞬时速度
        \item 求 ${t} = a\left( {a > 0}\right)$ 时的瞬时速度
        \item 借助(3)的结果，求 $t = \dfrac{5}{2}$ 时的瞬时速度
    \end{enumerate}
    \begin{jieda}
        \begin{enumerate}
            \item ${4.05g}{m}/{s},{4.005g}{m}/{s},{4.0005g}{m}/{s}$ \item ${4g}{m}/{s}$
            \item $a{g}{m}/{s}$
            \item $\dfrac{5}{2}g{m}/{s}$
        \end{enumerate}
    \end{jieda}
    
    \item 某水管的流水量 ${y}$ (单位: ${{m}}^{3}$ ) 与时间 ${t}$ (单位: ${s}$ ) 满足函数关系 ${y} = {f}\left( {t}\right)$ , 其中 ${f}\left( {t}\right)  = 3{t}$ .
    \begin{enumerate}
        \item 求 ${f}\left( {t}\right)$ 在 ${t} = {a}$ 处的导数 ${{f}}^{\prime }\left( {a}\right)$
        \item ${{f}}^{\prime }\left( {a}\right)$ 的实际意义是什么？
        \item 随着 ${a}$ 的变化 ${{f}}^{\prime }\left( {a}\right)$ 是否发生变化？为什么？
    \end{enumerate}
    \begin{jieda}
        \begin{enumerate}
            \item ${f}^{\prime }\left( a\right)  = 3$
            \item $t = a$ 时流水量的瞬时变化率
            \item 不变, ${f}^{\prime }\left( a\right)  = 3$ 为常数
        \end{enumerate}
    \end{jieda}
\end{enumerate}
\clearpage

\subsection*{导数的运算(1)初等函数的导数\&四则运算}
A 组：
\begin{enumerate}
    \item 初等函数的导数
    
    \begin{center}
    
    \begin{tabular}{|c|c|}
    \hline
    原函数 & 导函数 \\
    \cline{1-2}
    $f\left( x\right)  = c$ ( $c$ 为常数) & $f\left( x\right)  =$ \blkda{0}. \\
    \cline{1-2}
    $f\left( x\right)  = {x}^{\alpha }\left( {\alpha  \in  {Q},\text{ 且 }\alpha  \neq  0}\right)$ & $f\left( x\right)  =$ \blkda{$\alpha(x^{\alpha-1})$}. \\
    \cline{1-2}
    $f\left( x\right)  = \sin x$ & $f\left( x\right)  =$ \blkda{$cosx$}. \\
    \cline{1-2}
    $f\left( x\right)  = \cos x$ & $f\left( x\right)  =$ \blkda{$-sinx$}. \\
    \cline{1-2}
    $f\left( x\right)  = {a}^{x}\left( {a > 0,\text{ 且 }a \neq  1}\right)$ & $f\left( x\right)  =$ \blkda{$lna \cdot a^x$}. \\
    \cline{1-2}
    $f\left( x\right)  = {{e}}^{x}$ & $f\left( x\right)  =$ \blkda{$e^x$}. \\
    \cline{1-2}
    $f\left( x\right)  = {\log }_{a}x\left( {a > 0,\text{ 且 }a \neq  1}\right)$ & $f\left( x\right)  =$ \blkda{$\dfrac{1}{lna \cdot x}$}. \\
    \cline{1-2}
    $f\left( x\right)  = \ln x$ & $f\left( x\right)  =$ \blkda{$\dfrac{1}{x}$}. \\
    \cline{1-2}
    \hline
    \end{tabular}
    
    \end{center}

    \item 已知 $f\left( x\right) ,g\left( x\right)$ 为可导函数，且 $g\left( x\right)  \neq  0$ .
    \begin{enumerate}
        \item ${\left\lbrack  f\left( x\right)  \pm  g\left( x\right) \right\rbrack  }^{\prime } =$ \blkda[3]{$f'(x) \pm g'(x)$}.
        \item ${\left\lbrack  f\left( x\right)  \cdot  g\left( x\right) \right\rbrack  }^{\prime } =$ \blkda[3]{$f'(x)g(x)+g'(x)f(x)$}，特别地， ${\left\lbrack  cf\left( x\right) \right\rbrack  }^{\prime } =$ \blkda[3]{$cf'(x)$}.
        \item $\left[\dfrac{f(x)}{g\left( x\right)}\right]'=$\blkda[4]{$\dfrac{f'(x)g(x)-g'(x)f(x)}{g^2(x)}$}.
    \end{enumerate}
    
    \item 函数 $y = \ln x$ 在 $x = 2$ 处的导数是\blkda{$\dfrac{1}{2}$}.
    \begin{jieda}
        函数 $y = \ln x$ ,有 ${y}^{\prime } = \dfrac{1}{x}$ ,所以函数 $y = \ln x$ 在 $x = 2$ 处的导数是 $\dfrac{1}{2}$ .
    \end{jieda}
    
    \item 下列求导运算正确的是\dotfill{(\kh{D})}
    \xx{${\left( \sin 1\right) }^{\prime } = \cos 1$}{${\left( {x}^{3}\right) }^{\prime } = 3{x}^{4}$}{${\left( {{e}}^{x}\right) }^{\prime } = x{{e}}^{x - 1}$}{${\left( \cos x\right) }^{\prime } =  - \sin x$}
    \begin{jieda}
        ${\left( \sin 1\right) }^{\prime } = 0$ ,A 错误； 
        
        ${\left( {x}^{3}\right) }^{\prime } = 3{x}^{2}$ ,B 错误； 
        
        ${\left( {{e}}^{x}\right) }^{\prime } = {{e}}^{x}$ ,C 错误； 
        
        ${\left( \cos x\right) }^{\prime } =  - \sin x$, D正确
    \end{jieda}
    
    \item 求下列函数 $y = f\left( x\right)$ 的导数，其中:
    \begin{enumerate}
        \item $f\left( x\right)  = \sqrt[3]{{x}^{2}}$
        \item $f\left( x\right)  = {x}^{\pi }$ .
    \end{enumerate}
    \begin{jieda}
        \begin{enumerate}
            \item ${f}^{\prime }\left( x\right)  = \dfrac{2}{3}{x}^{-\frac{1}{3}}$
            \item ${f}^{\prime }\left( x\right)  = \pi {x}^{\pi  - 1}$
        \end{enumerate}
    \end{jieda}
    
    
    \item 函数 $f\left( x\right)  = \pi$ 的导 (函) 数是\blkda[3]{$f'(x) = 0$}.
    
    \item 函数 $f\left( x\right)  = \sqrt[{3}]{{x}^{5}}$ 的导 (函) 数是\blkda[3]{$f'(x) = \dfrac{5}{3}x^{\frac{2}{3}}$}.
    
    \item 函数 $f\left( x\right)  = \dfrac{1}{{x}^{3}}$ 的导(函)数是\blkda[3]{$f'(x) = -3x^{-4}$}.
\end{enumerate}

B组：
\begin{enumerate}
    \item 已知函数 $f\left( x\right)  = \sin x$ ，则 $\mathop{\lim }\limits_{{h \rightarrow  0}}\dfrac{f\left( {0 + h}\right)  - f\left( 0\right) }{h} =$ \blkda{1}.
    \begin{jieda}
        若 $f\left( x\right)  = \sin x$ ,则 ${f}^{\prime }\left( x\right)  = \cos x,\mathop{\lim }\limits_{{h \rightarrow  0}}\dfrac{f\left( {0 + h}\right)  - f\left( 0\right) }{h} = {f}^{\prime }\left( 0\right)  = \cos 0 = 1$ .
    \end{jieda}
    
    \item 若 $f\left( x\right)  = \sin x$ ，则 $\mathop{\lim }\limits_{{t \rightarrow  0}}\dfrac{f\left( {2t}\right)  - f\left( 0\right) }{t} =$ \blkda{2}.
    \begin{jieda}
        $\mathop{\lim }\limits_{{t \rightarrow  0}}\dfrac{f\left( {2t}\right)  - f\left( 0\right) }{t} = 2f'(0) = 2$
    \end{jieda}
    
    \item 已知 $f\left( x\right)  = {x}^{3} + \ln x + \sin x$ ，则 $f\left( x\right)$ 的导函数为\blkda[4]{$f'(x) = 3{x}^{2} + \dfrac{1}{x} + \cos x (x>  0)$}.
    
    \item 设函数 $f\left( x\right)$ 的导数为 ${f}^{\prime }\left( x\right)$ ,且 $f\left( x\right)  = {x}^{2} + {2x}{f}^{\prime }\left( 1\right)$ ,则 ${f}^{\prime }\left( 1\right)  =$ \blkda{$-2$}.
    \begin{jieda}
        $f'(x) = 2x+2f'(1)$，$f'(1) = 2+2f'(1)$，$f'(1) = -2$
    \end{jieda}
    
    \item 已知 $f\left( x\right)  = x \cdot  \sin x$ ，则 ${f}^{\prime }\left( x\right)  =$ \blkda[3]{$\sin x + x\cos x$}.
    \begin{jieda}
        ${f}^{\prime }\left( x\right)  = 1 \cdot  \sin x + x \cdot  \cos x = \sin x + x\cos x$ .
    \end{jieda}
    
    \item 某赛车启动时的位移s(米)和时间 $t$ (秒)的关系满足 $S\left( t\right)  = \dfrac{1}{6}{t}^{3} + 9{t}^{2} + {10t} - 5$ ， 则 $t = 4$ 时赛车的瞬时速度是\blkda{$90$}(米/秒).
    \begin{jieda}
        由导数的物理意义可得该质点在第 4 秒时的瞬时速度即函数 $S\left( t\right)$ 在 $t = 4$ 时的导数值，因为 $S\left( t\right)  = \dfrac{1}{6}{t}^{3} + 9{t}^{2} + {10t} - 5$，所以 ${S}^{\prime }\left( t\right)  = \dfrac{1}{2}{t}^{2} + {18t} + {10}$，所以 ${S}^{\prime }\left( 4\right)  = \dfrac{1}{2} \times  {16} + {18} \times  4 + {10} = {90}$，所以质点在第 3 秒时的瞬时速度为 90 (米/秒).
    \end{jieda}

    \item 无论我们对函数 $y = {{e}}^{x}$ 求多少次导数，结果仍然是它本身；这就像我们在生活中无论遇到多少艰难险阻，都要不忘初心，坚持自我，按照自己制定的目标，奋勇前行! 已知函数 $f\left( x\right)  = x \cdot  {{e}}^{x}$ ，则它的导函数 ${f}^{\prime }\left( x\right)  =$ \blkda{$\left( {1 + x}\right) {{e}}^{x}$}.
    \begin{jieda}
        根据导数的乘法运算法则，可知 ${\left( x \cdot  {{e}}^{x}\right) }^{\prime } = {x}^{\prime } \cdot  {{e}}^{x} + x \cdot  {\left( {{e}}^{x}\right) }^{\prime } = \left( {1 + x}\right) {{e}}^{x}$，所以, ${f}^{\prime }\left( x\right)  = \left( {1 + x}\right) {{e}}^{x}$ .
    \end{jieda}

    \item 设函数 $f\left( x\right)$ 的导函数为 ${f}^{\prime }\left( x\right)$ ，且 $f\left( x\right)  = \cos x - {f}^{\prime }\left( \dfrac{\pi }{6}\right) x$ ，则 $f\left( \dfrac{\pi }{2}\right)  =$ \blkda{$\dfrac{\pi}{8}$}.
    \begin{jieda}
        因为 ${f}^{\prime }\left( x\right)  =  - \sin x - {f}^{\prime }\left( \dfrac{\pi }{6}\right)$，所以 ${f}^{\prime }\left( \dfrac{\pi }{6}\right)  =  - \sin \dfrac{\pi }{6} - {f}^{\prime }\left( \dfrac{\pi }{6}\right)$ ,整理得 ${f}^{\prime }\left( \dfrac{\pi }{6}\right)  =  - \dfrac{1}{4}$，所以 $f\left( x\right)  = \cos x + \dfrac{1}{4}x$，故$f\left( \dfrac{\pi }{2}\right) = \dfrac{\pi}{8}$
    \end{jieda}

    \item 已知函数 $f\left( x\right)$ 的导数为 ${f}^{\prime }\left( x\right)$ ，若 $f\left( x\right)  = {f}^{\prime }\left( 1\right) {x}^{3} - {2}^{x}$ ，则 ${f}^{\prime }\left( 1\right)  =$ \blkda{$ln 2$}.
    \begin{jieda}
        因为 $f\left( x\right)  = {f}^{\prime }\left( 1\right) {x}^{3} - {2}^{x}$，则 ${f}^{\prime }\left( x\right)  = 3{f}^{\prime }\left( 1\right) {x}^{2} - {2}^{x}\ln 2$ ,令 $x = 1$ 可得 ${f}^{\prime }\left( 1\right)  = 3{f}^{\prime }\left( 1\right)  - 2\ln 2$，解得 ${f}^{\prime }\left( 1\right)  = \ln 2$ .
    \end{jieda}

    \item 若 $f\left( x\right)  = {x}^{2}\sin x$ ，则 ${f}^{\prime }\left( x\right)  =$ \blkda[3]{${2x}\sin x + {x}^{2}\cos x$}.
    \begin{jieda}
        ${f}^{\prime }\left( x\right)  = 2x \cdot sinx + x^2 \cdot cosx = {2x}\sin x + {x}^{2}\cos x$
    \end{jieda}

    \item 已知函数 ${f}_{1}\left( x\right)  = \sin x - \cos x,{f}_{2}\left( x\right)  = {f}_{1}^{\prime }\left( x\right) ,{f}_{3}\left( x\right)  = {f}_{2}\left( x\right) ,\cdots \cdots ,{f_n}\left( x\right)  = f_{n-1}{}^{\prime }\left( x\right)$ ，则 ${f}_{2020}\left( x\right)  =$ \blkda[3]{$-cosx-sinx$}.
    \begin{jieda}
        $f_2(x) = cosx+sinx$，$f_3(x) = -sinx+cosx$，$f_4(x) = -cosx-sinx$，$f_5(x) = sinx-cosx = f_1(x)$

        故$f_{2020}(x) = f_4(x) = -cosx-sinx$
    \end{jieda}

    \item 下列四组函数中，导数相等的是\dotfill{(\kh{D})}
    \xx{$f\left( x\right)  = 5$ 与 $g\left( x\right)  = {5x}$}{$f\left( x\right)  = 2{x}^{2} + 1$ 与 $g\left( x\right)  = {x}^{2} + 2$}{$f\left( x\right)  = 2 + \cos x$ 与 $g\left( x\right)  =  - \sin x$}{$f\left( x\right)  = 2\left( {1 - {x}^{2}}\right)$ 与 $g\left( x\right)  =  - 2{x}^{2} + 4$}
    
    \item 下列求导数运算正确的是\dotfill{(\kh{D})}
    \xx{${\left( {x}^{3}\right) }^{\prime } = 2{x}^{2}$}{${\left( \lg x\right) }^{\prime } = \dfrac{1}{x}$}{${\left( {x}^{3} + 5\right) }^{\prime } = 3{x}^{2} + 5$}{${\left( \sin x\cos x\right) }^{\prime } = \cos {2x}$}
    
    \item 已知某物体在平面上作变速直线运动，且位移 $s$ (单位:米)与时间 $t$ (单位: 秒) 之间的关系可用函数: $s = \ln \left( {t + 1}\right)  + {t}^{2} - t$ 表示,则该物体在 $t = 3$ 秒时的瞬时速度为\dotfill{(\kh{A})}
    \xx{$\dfrac{21}{4}$ 米/秒}{$\left( {6 + 2\ln 2}\right)$ 米/秒}{$\dfrac{21}{2}$ 米/秒}{$\left( {4 + \ln 2}\right)$ 米/秒}
    
    \item 给出定义: 设 ${f}^{\prime }\left( x\right)$ 是函数 $y = f\left( x\right)$ 的导函数, ${f}^{\prime \prime }\left( x\right)$ 是函数 $y = {f}^{\prime }\left( x\right)$ 的导函数,若方程 ${f}^{\prime \prime }\left( x\right)  = 0$ 有实数解 $x = {x}_{0}$ ,则称 $\left( {{x}_{0},f\left( {x}_{0}\right) }\right)$ 为函数 $y = f\left( x\right)$ 的 “拐点”. 经研究发现所有的三次函数 $f\left( x\right)  = a{x}^{3} + b{x}^{2} + {cx} + d\left( {a \neq  0}\right)$ 都有“拐点”,且该“拐点”也是函数 $y = f\left( x\right)$ 的图像的对称中心. 若函数 $f\left( x\right)  =  - {x}^{3} + 3{x}^{2}$ , 则 $f\left( \dfrac{1}{2021}\right)  + f\left( \dfrac{2}{2021}\right)  + f\left( \dfrac{3}{2021}\right)  + \cdots  + f\left( \dfrac{4040}{2021}\right)  + f\left( \dfrac{4041}{2021}\right)  =$\dotfill{(\kh{A})}
    \xx{8082}{-8082}{8084}{-8084}
    
    \item 求下列函数 $y = f\left( x\right)$ 的导数，其中:
    \begin{enumerate}
        \item $f\left( x\right)  = 3{e}^{x} - {x}^{e} + e$
        \item $f\left( x\right)  = \cos x - \dfrac{2}{x}$
        \item $f\left( x\right)  = {\left( 2x + 1\right) }^{3}$
        \item $f\left( x\right)  = \sqrt{x}\sin x$
        \item $f\left( x\right)  = x\ln x - \dfrac{1}{{x}^{2}}$
        \item $f\left( x\right)  = \dfrac{{x}^{2} - 1}{x}$
        \item $f\left( x\right)  = \dfrac{{x}^{2} - 1}{{x}^{2} + 1}$
        \item $f\left( x\right)  = \tan x$
    \end{enumerate}
    \begin{jieda}
        \begin{enumerate}
            \item $f'(x) = 3e^x - e x^{e-1}$
        
            \item $f'(x) = -\sin x + \dfrac{2}{x^2}$
        
            \item $f'(x) = 3(2x+1)^2 \cdot 2 = 6(2x+1)^2$
        
            \item $f'(x) = \dfrac{1}{2\sqrt{x}}\sin x + \sqrt{x}\cos x$
        
            \item $f'(x) = \ln x + 1 + \dfrac{2}{x^3}$
        
            \item $f'(x) = 1 + \dfrac{1}{x^2}$
        
            \item $f'(x) = \dfrac{(2x)(x^2+1) - (x^2-1)(2x)}{(x^2+1)^2} = \dfrac{4x}{(x^2+1)^2}$
        
            \item $f'(x) = \dfrac{\cos x \cdot \cos x - \sin x \cdot (-\sin x)}{\cos^2 x} = \dfrac{1}{\cos^2 x}$
        \end{enumerate}
    \end{jieda}
    
    \item 吹一个球形的气球时，气球半径 ${r}$ 将随着空气容量 ${V}$ 的增大而增大.
    \begin{enumerate}
        \item 写出气球半径 ${r}$ 关于气球内空气容量 ${V}$ 的函数表达式
        \item 求 ${V} = 1$ 时，气球的瞬时膨胀率(即气球半径关于气球内空气容量的瞬时变化率).
    \end{enumerate}
    \begin{jieda}
        \begin{enumerate}
            \item $r = \sqrt[3]{\dfrac{3V}{4\pi }}\left( {V \geq  0}\right)$
            \item ${\left. {r}^{\prime }\right| }_{V = 1} = \dfrac{1}{\sqrt[3]{36\pi }}$
        \end{enumerate}
    \end{jieda}
    
    \item 蜥蜴的体温与阳光照射的关系近似满足函数关系式: $T\left( t\right)  = \dfrac{120}{t + 5} + {15}$ ，其中 $T\left( t\right)$ 为蜥蜴的体温(单位: ${}^{ \circ  }{C}$ )， $t$ 为太阳落山后的时间(单位: ${{min}}$ )
    \begin{enumerate}
        \item 求 ${T}^{\prime }\left( {10}\right)$ ，并解释其实际意义
        \item 蜥蜴体温的瞬时变化率为 $- {1}^{ \circ  }{C}$ /min 时的时刻 $t$ 是多少(精确到 0.01)？
    \end{enumerate}
    \begin{jieda}
        \begin{enumerate}
            \item ${T}^{\prime }\left( t\right)  = \dfrac{-{120}}{{\left( t + 5\right) }^{2}}$ ,则 ${T}^{\prime }\left( {10}\right)  = \dfrac{-{120}}{{\left( {10} + 5\right) }^{2}} =  - \dfrac{8}{15}$，表示太阳落山后 ${10}{{min}}$ ,蜥蜴的体温下降的速度为 $- \dfrac{8}{15}{}^{ \circ  }{C}/{{min}}$ .
            
            \item 令 ${T}^{\prime }\left( t\right)  = \dfrac{-{120}}{{\left( t + 5\right) }^{2}} =  - 1$ ,解得 $t = 2\sqrt{30} - 5 \approx  2 \times  {5.477} - 5 \approx  {5.95}$，故蜥蜴体温的瞬时变化率为 $- {1}^{ \circ  }{C}/{{min}}$ 时的时刻是 ${5.95}{{min}}$ .
        \end{enumerate}
    \end{jieda}
\end{enumerate}

\clearpage
\subsection*{导数的概念和意义(2)导数的几何意义}
A 组：

\begin{enumerate}
    \item 已知 $f\left( x\right)  = a{x}^{3} + 3{x}^{2} + 2$ ，且 ${f}^{\prime }\left( {-1}\right)  = 4$ ，则实数 $a$ 的值为\blkda{$\dfrac{10}{3}$}.
    \begin{jieda}
        $f'(-1) = 3a-6 = -4$，故$a = \dfrac{10}{3}$
    \end{jieda}

    \item 若函数 $f\left( x\right)$ 在 $x = 3$ 处的导数 ${f}^{\prime }\left( 3\right)  = \sqrt{3}$ ，则曲线 $f\left( x\right)$ 在 $\left( {3,f\left( 3\right) }\right)$ 处的切线的倾斜角 $\theta  =$ \blkda{$\ang{60}$}.
    \begin{jieda}
        因为函数 $f\left( x\right)$ 在 $x = 3$ 处的导数 ${f}^{\prime }\left( 3\right)  = \sqrt{3}$，所以函数 $f\left( x\right)$ 在点 $\left( {3,f\left( 3\right) }\right)$ 处的切线斜率 $k = {f}^{\prime }\left( 3\right)  = \sqrt{3}$，所以 $\tan \theta  = \sqrt{3}$ ,又 ${0}^{ \circ  } \leq  \theta  < {180}^{ \circ  }$ ,所以倾斜角 $\theta  = {60}^{ \circ  }$ .
    \end{jieda}

    \item 已知函数 ${y} = {f}\left( {x}\right)$ 在 ${x} = {1}$ 处的切线方程为 ${y} = {4}{x} - {3}$ ，求 ${f}\left( {1}\right)$ 和 ${{f}}^{\prime }\left( {1}\right)$ .
    \begin{jieda}
        $f(1) = 1, f'(1) = 4$
    \end{jieda}

    \item 已知 ${f}\left( {x}\right)  = 3{{x}}^{2}$ ,分别求曲线 ${y} = {f}\left( {x}\right)$ 在点 ${P}\left( {-1,3}\right)$ 和点 ${Q}\left( {1,3}\right)$ 处的切线方程.
    \begin{jieda}
        $f(x) = 6x$，$f'(1) = 6, f'(-1) = -6$，故$P$处切线为$y = -6x-3$，$Q$处切线为$y = 6x-3$.
    \end{jieda}


    \item 如图，已知直线 ${l}$ 是曲线 ${y} = {f}\left( {x}\right)$ 在 ${x} = {3}$ 处的切线,求 ${{f}}^{\prime }\left( {3}\right)$ .
    \begin{center}
    \includegraphics[width=0.3\textwidth]{bodyxb/src/bo_d6fgfrc601uc73e0627g_0_215_1016_411_283_0.jpg}
    \end{center}
    \begin{jieda}
        根据导数的几何意义可知$f'(3) = -\dfrac{1}{3}$
    \end{jieda}

\end{enumerate}

B 组：
\begin{enumerate}
    \item 函数 $f\left( x\right)  = \dfrac{\sqrt{3}}{3}{x}^{3} - 1$ 在 $x = 1$ 处的切线倾斜角是\blkda{$\ang{60}$}.
    \begin{jieda}
        $f'(1) = \sqrt{3}$，故该处切线倾斜角为$\ang{60}$.
    \end{jieda}

    \item 抛物线 $\Gamma  : y = {x}^{2}$ 上一点 $M$ 到焦点的距离为 1,抛物线 $\Gamma$ 在点 $M$ 处的切线的斜率为\blkda{$\pm \sqrt{3}$}.
    \begin{jieda}
        $\Gamma  : y = {x}^{2}$的准线为$y = -\dfrac{1}{4}$，故$M$的坐标为$\left(\pm \dfrac{\sqrt{3}}{2}, \dfrac{3}{4}\right)$，$f'\left(\pm \dfrac{\sqrt{3}}{2}\right) = \pm \sqrt{3}$
    \end{jieda}

    \item 求曲线 $y = \cos x$ 在 $x = \dfrac{\pi }{2}$ 处的切线方程是\blkda[3]{$y = -\left(x-\dfrac{\pi}{2}\right)$}.
    \begin{jieda}
        设$f(x) = cosx$，$f'(x) = -sinx$，$f'\left(\dfrac{\pi}{2}\right) = -1$，故切线为$y = -\left(x-\dfrac{\pi}{2}\right)$
    \end{jieda}

    \item 求过点 $\left( {0, - 1}\right)$ 且与曲线 $y = 2{x}^{2}$ 相切的直线的方程是\blkda[3]{$y = \pm 2\sqrt{2}x-1$}.
    \begin{jieda}
        设切点为$(m, 2m^2)$，该处切线为$y = 4m(x-m)+2m^2$，令其过点 $\left( {0, - 1}\right)$，得到$m = \pm \dfrac{\sqrt{2}}{2}$，故切线为$y = \pm 2\sqrt{2}x-1$
    \end{jieda}

    \item 过点 $\left( {1,2}\right)$ 且与 $y = 2{x}^{2}$ 相切的直线方程为\blkda[3]{$y = 4x-2$}.

    \item $y = {kx} + b$ 是 $y = \dfrac{\ln x}{{x}^{2}}$ 在 $\left( {1,0}\right)$ 处的切线方程,则 $k - b =$\blkda{$2$}
    \begin{jieda}
        ${y}^{\prime } = \dfrac{\dfrac{1}{x} \cdot  {x}^{2} - {2x}\ln x}{{x}^{4}} = \dfrac{1 - 2\ln x}{{x}^{3}},x > 0,y'|_{x=1}  = 1$ ,且$x=1$时$y = 0$，所以 $\left\{  \begin{array}{l} k = 1 \\  k + b = 0 \end{array}\right.$ ,得 $k = 1,b =  - 1$ ,则 $k - b = 2$ .
    \end{jieda}

    \item 曲线 $y = \ln x - \dfrac{2}{x}$ 在 $x = 1$ 处的切线方程为\blkda[3]{${3x} - y - 5 = 0$}.
    \begin{jieda}
        设 $f\left( x\right)  = \ln x - \dfrac{2}{x},x > 0$ ,则 ${f}^{\prime }\left( x\right)  = \dfrac{1}{x} - 2\left( {-\dfrac{1}{{x}^{2}}}\right)  = \dfrac{1}{x} + \dfrac{2}{{x}^{2}}$，所以 ${f}^{\prime }\left( 1\right)  = 3$ ,且 $f\left( 1\right)  =  - 2$ ,即直线斜率 $k = 3$ ,过点$(1, -2)$，故曲线 $y = \ln x - \dfrac{2}{x}$ 在 $x = 1$ 处的切线方程为 $y + 2 = 3\left( {x - 1}\right)$ ,即 ${3x} - y - 5 = 0$
    \end{jieda}

    \item 已知 $f\left( x\right)$ 为 ${R}$ 上的奇函数，当 $x < 0$ 时， $f\left( x\right)  =  - {x}^{2} + x$ ，则 $f\left( x\right)$ 的图象在 $x = 1$ 处的切线方程为\blkda[3]{${3x} - y - 1 = 0$}.
    \begin{jieda}
        设 $x > 0$ ,则 $- x < 0,f\left( {-x}\right)  =  - {x}^{2} - x$ ,又 $f\left( x\right)$ 为 $\mathbf{R}$ 上的奇函数，所以 $f\left( x\right)  =  - f\left( {-x}\right)  = {x}^{2} + x$，即当 $x > 0$ 时, $f\left( x\right)  = {x}^{2} + x$ ,当 $x > 0$ 时, ${f}^{\prime }\left( x\right)  = {2x} + 1,{f}^{\prime }\left( 1\right)  = 3,f\left( 1\right)  = 2$，所以 $f\left( x\right)$ 的图象在 $x = 1$ 处的切线方程为 $y - 2 = 3\left( {x - 1}\right)$ ,即 ${3x} - y - 1 = 0$ .
    \end{jieda}

    \item 若点 $P$ 是曲线 $y = {x}^{2} - \ln x$ 上的任意一点，则点 $P$ 到直线 $y = x - 2$ 的最小距离是\blkda{$\sqrt{2}$}.
    \begin{jieda}
        设斜率为 1 的直线 $l$ 与曲线 $y = {x}^{2} - \ln x$ 相切于点 $Q\left( {{x}_{0},{x}_{0}^{2} - \ln {x}_{0}}\right) \left( {{x}_{0} > 0}\right)$ 
        
        因为 $y = {x}^{2} - \ln x\left( {x > 0}\right)$ ,所以 ${y}^{\prime } = {2x} - \dfrac{1}{x}$
        
        令 ${y}^{\prime } = {2x} - \dfrac{1}{x} = 1$ ,解得 $x = 1$ ,所以 $Q\left( {1,1}\right)$，所以切线 $l$ 的方程为: $y = x$，所以要求点 $P$ 到直线 $y = x - 2$ 的最小距离，即求切线 $l : y = x$ 到直线 $y = x - 2$ 的距离，由两平行线间的距离公式可得 $d = \dfrac{2}{\sqrt{{1}^{2} + {1}^{2}}} = \sqrt{2}$，所以点 $P$ 到直线 $y = x - 2$ 的最小距离是 $\sqrt{2}$ .
    \end{jieda}

    \item 曲线 $f\left( x\right)  =  - {2x} + \ln x$ 在点 $\left( {1,f\left( 1\right) }\right)$ 处的切线与两坐标轴所围成的三角形的面积为\blkda{$\dfrac{1}{2}$}.
    \begin{jieda}
        已知 $f\left( x\right)  =  - {2x} + \ln x$ ,得： ${f}^{\prime }\left( x\right)  =  - 2 + \dfrac{1}{x}$ ,所以 ${f}^{\prime }\left( 1\right)  =  - 1$ ,又 $f\left( 1\right)  =  - 2$，所以曲线 $f\left( x\right)  =  - {2x} + \ln x$ 在点 $\left( {1,f\left( 1\right) }\right)$ 处的切线方程为 $y - \left( {-2}\right)  =  - \left( {x - 1}\right)$ ,整理得$y =  - x - 1$ .
        
        直线 $y =  - x - 1$ 与 $x$ 轴交于点(-1,0),与 $y$ 轴交于点$(0, -1)$，因此所求三角形的面积为 $\dfrac{1}{2} \times  1 \times  1 = \dfrac{1}{2}$ 
    \end{jieda} 

    \item 函数特性 $P$ : “函数的图像上存在两点,使得函数的图像在这两点处的切线 互相垂直”，则下列函数中满足特性 $P$ 的函数为\dotfill{(\kh{B})}
    \xx{$y = {x}^{3}$}{$y = \sin x$}{$y = {{e}}^{x}$}{$y = \ln x$}
    \begin{jieda}
        问题转化为导函数的值域中存在两个元素之积为$-1$.
    \end{jieda}
    
    \item 过点 $\left( {2, - 6}\right)$ 作曲线 $y = {x}^{3} - {3x}$ 的切线,所得切线斜率为\dotfill{(\kh{C})}
    \xx{$-3$}{0 或 3}{$-3$ 或 24}{0}
    
    \item 求曲线 $y = {x}^{3} - {3x} + 5$ 平行于直线 $x$ 轴的切线及其切点坐标.
    \begin{jieda}
        $y = 3$ ,切点 $\left( {1,3}\right)$
        
        $y = 7$ ,切点$(-1,7)$
    \end{jieda}
    
    \item 求曲线 $y = \dfrac{1}{x}$ 平行于直线 $y =  - x$ 的切线及其切点坐标.
    \begin{jieda}
        $y =  - x + 2$ ,切点(1,1)； $y =  - x - 2$ ,切点$(-1, -1)$
    \end{jieda}
    
    \item 直线 $y =  - x + b$ 是下列函数的切线吗? 如果是,请求出 ${b}$ 的值; 如果不是，请说明理由. \begin{enumerate*}
        \item $y = \ln x$
        \item $y = \dfrac{1}{x}$
    \end{enumerate*} 
    \begin{jieda}
        \begin{enumerate}
            \item 不是, ${y}^{\prime } = \dfrac{1}{x} \in  \left( {0, + \infty }\right)$
            \item $b =  \pm  2$
        \end{enumerate}
    \end{jieda}
    
    \item 已知函数 $f\left( x\right)  = \sin x - x + a{x}^{2},a \in  {R}$ .
    \begin{enumerate}
        \item 若曲线 $y = f\left( x\right)$ 在 $x = \pi$ 处的切线过原点，求 $a$ 的值
        \item 当 $x \leq  5$ 时， $f\left( x\right)  \geq  0$ ，求 $a$ 的取值范围
    \end{enumerate}
    \begin{jieda}
        \begin{enumerate}
            \item $f\left( x\right)  = \sin x - x + {ax^2}$ ,则 ${f}^{\prime }\left( x\right)  = \cos x - 1 + {2ax}$，则 $k = {f}^{\prime }\left( \pi \right)  = \cos \pi  - 1 + {2a\pi } =  - 2 + {2a\pi }$
            
            又曲线 $y = f\left( x\right)$ 在 $x = \pi$ 处的切线过原点，则 $\dfrac{f\left( \pi \right)  - 0}{\pi  - 0} = k$ ,即 $- 1 + {a\pi } =  - 2 + {2a\pi }$ ,解之得 $a = \dfrac{1}{\pi }$
            
            \item 
            \begin{enumerate}
                \item 方法1

                $f(0) = 0$，$f\left(\pi\right) = a\pi^2-\pi \geq 0$，故$a \geq \dfrac{1}{\pi}$
    
                设$g(x) = \sin x - x + \dfrac{x^2}{\pi}$，故$f(x) \geq g(x)$，$g(0) = g(\pi) = 0$，$g'(x) = \dfrac{2x}{\pi}+cosx-1$
                
                \begin{dingautolist}{172}
                    \item 当$x < 0$时，$g'(x) < 0$，故$g(x)$在$(-\infty, 0)$严减，由$g(0) = 0$可知$g(x) > 0$
                    \item 当$x \in (0, \pi)$
    
                    $g'\left(\dfrac{\pi}{2}\right) = 0$，$g'\left(\pi\right) = 0$
                    
                    $g''(x) = \dfrac{2}{\pi} - sinx$，故$g''(x)$在存在变号零点$x_1 = arcsin\left(\dfrac{2}{\pi}\right), x_2 = \pi - arcsin\left(\dfrac{2}{\pi}\right)$
    
                    故$g'(x)$在$(0, x_1)$严增，在$(x_1, x_2)$严减，在$(x_2, \pi)$严增
    
                    因此$g(x)$在$\left(0, \dfrac{\pi}{2}\right)$严增，在$\left(\dfrac{\pi}{2}, \pi\right)$严减，故$g(x) > 0$
    
                    \item 当$x \in (\pi, 5]$，$g'(x) \geq 0$恒成立，$g(x)$严增，$g(x) \geq 0$
                \end{dingautolist}
    
                综上,当 $a \geq  \dfrac{1}{\pi }$ 时,  $f\left( x\right) \geq g(x) \geq  0$ 在 $( - \infty ,5\rbrack$ 上恒成立.

                \item 方法2
                
                $f(0) = 0$，满足要求，故当$x \neq 0$时，应有$a \geq \dfrac{x-sinx}{x^2}$在$(-\infty, 0)\cup (0, 5]$恒成立

                设$g(x) = \dfrac{x-sinx}{x^2}$，$g'(x) = \dfrac{(1-cosx)x^2-2x(x-sinx)}{x^4} = \dfrac{x(1-cosx)-2(x-sinx)}{x^3} = \dfrac{2sinx - x(cosx+1)}{x^3} = \dfrac{4sin\left(\dfrac{x}{2}\right)cos\left(\dfrac{x}{2}\right)-x\left(2cos^2\left(\dfrac{x}{2}\right)\right)}{x^3} = \dfrac{2cos\left(\dfrac{x}{2}\right)\left[2sin\left(\dfrac{x}{2}\right)-xcos\left(\dfrac{x}{2}\right)\right]}{x^3}$

                令$g'(x) = 0$

                \begin{dingautolist}{172}
                    \item $cos\left(\dfrac{x}{2}\right) = 0$，$x = \pi + 2k\pi, k \in \mathbf{Z}$，此时$sinx = 0$，故$g(x) = \dfrac{1}{x} \leq \dfrac{1}{\pi}$ （$x = \pi$处取等）

                    在$\pi$左侧，$g'(x) > 0$，在$\pi$右侧，$g'(x) < 0$，故$g(\pi)$为$g(x)$的极大值.

                    \item $2sin\left(\dfrac{x}{2}\right)-xcos\left(\dfrac{x}{2}\right) = 0$，即$tan\left(\dfrac{x}{2}\right) = \dfrac{x}{2}$，此时$g(x) = \dfrac{2tan\left(\dfrac{x}{2}\right)-\dfrac{2tan\left(\dfrac{x}{2}\right)}{1+tan^2\left(\dfrac{x}{2}\right)}}{\left(2tan\left(\dfrac{x}{2}\right)\right)^2} = \dfrac{1-\dfrac{1}{1+tan^2\left(\dfrac{x}{2}\right)}}{2tan\left(\dfrac{x}{2}\right)} = \dfrac{tan\left(\dfrac{x}{2}\right)}{2\left(1+tan^2\left(\dfrac{x}{2}\right)\right)} \leq \dfrac{1}{4}$
                \end{dingautolist}
                综上，$g(\pi) = \dfrac{1}{\pi}$为$g(x)$的最大值，故当 $a \geq  \dfrac{1}{\pi }$ 时,  $f\left( x\right) \geq g(x) \geq  0$ 在 $( - \infty ,5\rbrack$ 上恒成立.
            \end{enumerate}    
            \item 方法3

            $f(x) \geq 0$即$-ax^2 + x\leq sinx$
            \begin{dingautolist}{172}
                \item 当$x = 0$时，显然成立
                \item 当$x > 0$时，$-ax + 1 \leq \dfrac{sinx}{x}$，即过$(0, 1)$的直线在$y = \dfrac{sinx}{x}$图像及其下方，计算相切情况可知$a \geq \dfrac{1}{\pi}$
                \item 当$x < 0$时，$-ax + 1 \geq \dfrac{sinx}{x}$，即过$(0, 1)$的直线在$y = \dfrac{sinx}{x}$图像及其上方，计算可知$a \geq 0$
            \end{dingautolist}
            综上，$a \geq \dfrac{1}{\pi}$
        \end{enumerate}
    \end{jieda}
    
    \item 已知曲线 ${C}_{1} : y = {x}^{3}$ 和 ${C}_{2} : y = a{x}^{2} + x - 2,\left( {a \in  {R}}\right)$ .
    \begin{enumerate}
        \item 若曲线 ${C}_{1}$ 、 ${C}_{2}$ 在 $x = 1$ 处的切线互相垂直，求 $a$ 的值
        \item 若存在$x_0$，使得曲线 ${C}_{1}$ 、 ${C}_{2}$ 在 $x = {x}_{0}$ 处的切线平行或重合且斜率大于 3，求 $a$ 的取值范围.
    \end{enumerate}
    \begin{jieda}
        \begin{enumerate}
            \item 由 $y = {x}^{3}$ 可得 ${y}^{\prime } = 3{x}^{2}$ ,
            由 $y = a{x}^{2} + x - 2,\left( {a \in  \mathbf{R}}\right)$ 可得 ${y}^{\prime } = {2ax} + 1$ ,
            因为曲线 ${C}_{1}\text{、}{C}_{2}$ 在 $x = 1$ 处的切线互相垂直,
            所以 ${k}_{1} \cdot  {k}_{2} = \left( {3 \times  {1}^{2}}\right)  \times  \left( {{2a} + 1}\right)  =  - 1$ ,解得 $a =  - \dfrac{2}{3}$ .
            \item 由题意,切线的斜率 $k = 3{x}_{0}^{2} = {2a}{x}_{0} + 1 > 3$ ,
            可得 $3{x}_{0}^{2} - 1 = {2a}{x}_{0}$ ,且 ${x}_{0} > 1$ 或 ${x}_{0} <  - 1$ 
            
            所以 ${2a} = 3{x}_{0} - \dfrac{1}{{x}_{0}}$ ,
            令 $h\left( x\right)  = {3x} - \dfrac{1}{x}$ ,则函数在 $\left( {1, + \infty }\right)$ 和 $\left( {-\infty , - 1}\right)$ 上是增函数,
            所以 $h\left( x\right)  > h\left( 1\right)  = 2$ 或 $h\left( x\right)  < h\left( {-1}\right)  =  - 2$ ,
            即 ${2a} > 2$ 或 ${2a} <  - 2$ ,解得 $a > 1$ 或 $a <  - 1$ .
        \end{enumerate}
    \end{jieda}
\end{enumerate}
\clearpage
\subsection*{导数的运算(2)简单复合函数的导数}
A组：
\begin{enumerate}
    \item 求下列函数的导数:
    \begin{enumerate*}
        \item $y = {\left( 3 - 2x\right) }^{2}$
        \item $y = \sin {2x}$
    \end{enumerate*}
    \begin{jieda}
        \begin{enumerate}
            \item $y' = 2(3-2x)\cdot (-2) = 8x-12$
            \item $y' = cos2x \cdot 2 = 2cos2x$
        \end{enumerate}
    \end{jieda}
    
    \item 下列求导运算正确的是\dotfill{(\kh{A})}
    \xx{${\left( \ln x\right) }^{\prime } = \dfrac{1}{x}$}{${\left\lbrack  \ln \left( 2x - 1\right) \right\rbrack  }^{\prime } = \dfrac{1}{{2x} - 1}$}{${\left( {2}^{x}\right) }^{\prime } = {2}^{x}$}{${\left( \sin x\right) }^{\prime } =  - \cos x$}
    \begin{jieda}
        对于 ${A}$ 选项,由对数函数的求导公式,得 ${\left( \ln x\right) }^{\prime } = \dfrac{1}{x}$ ,故 ${A}$ 正确；
        
        对于 ${B}$ 选项,由复合函数的求导法则,得 ${\left\lbrack  \ln \left( 2x - 1\right) \right\rbrack  }^{\prime } = \dfrac{1}{{2x} - 1} \cdot  {\left( 2x - 1\right) }^{\prime } = \dfrac{2}{{2x} - 1}$ ,故 ${B}$ 错误;
        
        对于 ${C}$ 选项,由指数函数的求导公式,得 ${\left( {2}^{x}\right) }^{\prime } = {2}^{x}\ln 2$ ,故 ${C}$ 错误；
        
        对于 ${D}$ 选项,由正弦函数的求导公式,得 ${\left( \sin x\right) }^{\prime } = \cos x$ ,故 ${D}$ 错误.
    \end{jieda}
\end{enumerate}
B组：
\begin{enumerate}
    \item 已知函数 $f\left( x\right)  = {{e}}^{{2x} + 1}$ ，则 $\mathop{\lim }\limits_{{{\Delta x} \rightarrow  0}}\dfrac{f\left( {1 - {\Delta x}}\right)  - f\left( 1\right) }{\Delta x} =$ \blkda{$- 2{{e}}^{3}$}.
    \begin{jieda}
        由 $\mathop{\lim }\limits_{{{\Delta x} \rightarrow  0}}\dfrac{f\left( {1 - {\Delta x}}\right)  - f\left( 1\right) }{\Delta x} =  - \mathop{\lim }\limits_{{{\Delta x} \rightarrow  0}}\dfrac{f\left( 1\right)  - f\left( {1 - {\Delta x}}\right) }{\Delta x} =  - {f}^{\prime }\left( 1\right)$ 
        
        又由函数 $f\left( x\right)  = {{e}}^{{2x} + 1}$ ,可得 ${f}^{\prime }\left( x\right)  = 2{{e}}^{{2x} + 1}$ ,则 ${f}^{\prime }\left( 1\right)  = 2{{e}}^{3}$
        
        所以 $\mathop{\lim }\limits_{{{\Delta x} \rightarrow  0}}\dfrac{f\left( {1 - {\Delta x}}\right)  - f\left( 1\right) }{\Delta x} =  - 2{{e}}^{3}$ .
    \end{jieda}
    
    \item 曲线 $f\left( x\right)  = \dfrac{{{e}}^{x} - {x}^{2}}{x} + x\ln \left( {{e}x}\right)$ 在点 $\left( {1,f\left( 1\right) }\right)$ 处的切线方程为\blkda[3]{$x - y + {e} - 1 = 0$}.
    \begin{jieda}
        因为 $f\left( x\right)  = \dfrac{{{e}}^{x} - {x}^{2}}{x} + x\ln \left( {{e}x}\right)  = \dfrac{{{e}}^{x}}{x} + x \cdot  \ln x$ ,所以$f\left( 1\right)  = {e},{f}^{\prime }\left( x\right)  = \dfrac{\left( {x - 1}\right) {{e}}^{x}}{{x}^{2}} + 1 + \ln x,{f}^{\prime }\left( 1\right)  = 1$，所以切线方程为 $y - {e} = x - 1$ ,即 $x - y + {e} - 1 = 0$ .
    \end{jieda}
    
    \item 曲线 $y = {2}^{1 - {3x}}$ 在点 $\left( {0,2}\right)$ 处的切线方程为\blkda[3]{$y = (-6ln2)x+2$}.
    \begin{jieda}
        $y' = ln2(2^{1-3x}) \cdot (-3)$，$y'|_{x=0} = -6ln2$，故切线方程为$y = (-6ln2)x+2$
    \end{jieda}
    
    \item 已知曲线 $y = a\ln {3x}\left( {a > 0}\right)$ 在 $x = 1$ 处的切线与两坐标轴围成的三角形的面积为 $4{\left( \ln 3 - 1\right) }^{2}$ ，则 $a =$ \blkda{8}.
    \begin{jieda}
        $y' = \dfrac{a}{x}$，故切线为$y = a(x-1)+aln3$，切线和两坐标轴交点为$(0, a(ln3-1)), (1-ln3, 0)$，又因为$a>0$，故$S = \dfrac{a(ln3-1)^2}{2}$，因此$a = 8$
    \end{jieda}
    
    \item 已知奇函数 $y = f\left( x\right)$ 的定义域为 $\mathbf{R}$ ,当 $x < 0$ 时, $f\left( x\right)  = {x}^{3} + {{e}}^{x} - 1$ ,则 ${f}^{\prime }\left( 1\right)  =$ \blkda{$3+e^{-1}$}.
    \begin{jieda}
        $f'(1) = f'(-1) = 3(-1)^2+e^{-1} = 3+e^{-1}$
    \end{jieda}
    
    \item 设 $f\left( x\right)  = \cos \left( {\sqrt{3}x + \varphi }\right) \left( {0 < \varphi  < \pi }\right)$ ，若 $y = f\left( x\right)  + {f}^{\prime }\left( x\right)$ 是奇函数，则 $\varphi  =$\blkda{$\dfrac{\pi}{6}$}
    \begin{jieda}
        $f'(x) = -sin(\sqrt{3}x+\varphi) \cdot (\sqrt{3})$

        $y = f(x) + f'(x) = 2sin\left(\sqrt{3}x+\varphi+\dfrac{5\pi}{6}\right)$，因为是奇函数，所以$\varphi+\dfrac{5\pi}{6} = k\pi, k \in \mathbf{Z}$，故$\varphi = \dfrac{\pi}{6}$
    \end{jieda}
    
    \item 求下列函数的导数:
    \begin{enumerate}
        \item $f\left( x\right)  = \dfrac{x - 1}{x - 2}$ ，则 ${f}^{\prime }\left( x\right)  =$ \blkda{$- \dfrac{1}{{\left( x - 2\right) }^{2}}$}.
        \item $f\left( x\right)  = 3{x}\sqrt{2 - x}$ ，则 ${f}^{\prime }\left( x\right)  =$ \blkda[3]{$-\dfrac{3x}{2\sqrt{2 - x}} + 3\sqrt{2 - x}$}.
        \item $f\left( x\right)  = \dfrac{\ln \left( {{2x} + 1}\right) }{x}$ ，则 ${f}^{\prime }\left( x\right)  =$ \blkda[4]{$\dfrac{2}{2{x}^{2} + x} - \dfrac{\ln \left( {{2x} + 1}\right) }{{x}^{2}}$}.
        \item $f\left( x\right)  = {e}^{x}\cos {2x}$ ，则 ${f}^{\prime }\left( x\right)  =$ \blkda[3]{${e}^{x}(\cos 2x - 2\sin 2x)$}.
    \end{enumerate}
    
    \item 已知函数 $f\left( x\right)  = \sqrt{3}\sin \left( {{\omega x} - \dfrac{\pi }{6}}\right) \left( {\omega  > 0}\right)$ 的图象向左平移 $\varphi \left( {\varphi  > 0}\right)$ 个单位长度后与其导函数 $y = {f}^{\prime }\left( x\right)$ 的图像重合,则 $f\left( \varphi \right)$ 的值为\blkda{$\dfrac{3}{2}$}.
    \begin{jieda}
        $f'(x) = \sqrt{3}\omega cos\left( {{\omega x} - \dfrac{\pi }{6}}\right)$

        由图像重合可知$\omega = 1$，因此$f(x) = \sqrt{3}sin\left(x-\dfrac{\pi}{6}\right)$，$f'(x) = \sqrt{3}cos\left( {x - \dfrac{\pi }{6}}\right)$

        故根据诱导公式可知$\varphi = \dfrac{\pi}{2}$，因此$f(\varphi) = \dfrac{3}{2}$
    \end{jieda}
    
    \item 一个球的半径以 2 厘米/秒的速度在增长，当此球的半径为 6 厘米时，其体积的变化率为\dotfill{(\kh{B})}
    \xx{$576\pi {cm}^{3}/$ 秒}{$288\pi {cm}^{3}/$ 秒}{$144\pi {cm}^{3}/$ 秒}{${96\pi } {cm}^{3}/$ 秒}
    
    \item 一个圆的半径以 0.1 厘米/秒的速度在增长，若圆的周长为 $C$ 厘米，则其面积的变化率为\dotfill{(\kh{B})}
    \xx{$0.2\pi C {cm}^{2}/$ 秒}{$0.1C {cm}^{2}/$ 秒}{$\dfrac{0.1C}{2\pi} {cm}^{2}/$ 秒}{$-\dfrac{0.1C}{2\pi} {cm}^{2}/$ 秒}
    
    \item 已知函数 $y = f\left( x\right)$ 及其导函数 $y = {f}^{\prime }\left( x\right)$ 的定义域均为$\mathbf{R}$,记 $g\left( x\right)  = {f}^{\prime }\left( x\right)$ . 若 $y = f\left( {\dfrac{3}{2} - {2x}}\right) \text{ 、 }y = g\left( {2 + x}\right)$ 均为偶函数,有下列两个结论:  \ding{172}  $f\left( {-1}\right)  = f\left( 4\right)$ ;  \ding{173}  $g\left( {-\dfrac{1}{2}}\right)  = 0$ ，则\dotfill{(\kh{C})}
    \xx{\ding{172} 对 \ding{173} 错}{\ding{172} 错 \ding{173} 对}{\ding{172} 对 \ding{173} 对}{\ding{172} 错 \ding{173} 错}
    \begin{jieda}
        $f(x)$关于$x = \dfrac{3}{2}$对称，故\ding{172}正确，$g(x)$关于$x = 2$对称

        因为两函数定义域均为$\mathbf{R}$，故$g(x)$有对称中心$\left(\dfrac{3}{2}, 0\right)$，故$g(x)$周期为$2$，故\ding{173}正确.
    \end{jieda}
    
    \item 求下列函数的导数:
    \begin{enumerate}
        \item ${y} = \dfrac{1}{\sqrt{1 - 2{x}}}$
        \item $y = {\log }_{2}\left( {{2x} + 1}\right)$
        \item $y = {e}^{{3x} + 2}$
        \item $y = \sin \left( {{2x} + \dfrac{\pi }{3}}\right)$
        \item $y = {e}^{2x+ 1}\sin x$
        \item $f\left( x\right)  = {\left( -x + 5\right) }^{\frac{5}{3}}$
        \item $f\left( x\right)  = \dfrac{\ln \left( {3x}\right) }{{2x} + 1}$
    \end{enumerate}
    \begin{jieda}
        \begin{enumerate}
            \item $y = {\left( 1 - 2x\right) }^{-\frac{1}{2}}$ ,设 $y = {u}^{-\frac{1}{2}},u = 1 - {2x}$，则 ${y}^{\prime }x = {y}^{\prime }u \cdot  {u}^{\prime }x = {\left( {u}^{-\frac{1}{2}}\right) }^{\prime }{\left( 1 - 2x\right) }^{\prime } = \left( {-\dfrac{1}{2}{{u}}^{-\frac{3}{2}}}\right)  \cdot  \left( {-2}\right)  = {\left( 1 - 2{x}\right) }^{-\frac{3}{2}}.$即 ${{y}}^{\prime } = {\left( 1 - 2{x}\right) }^{-\frac{3}{2}}$ .
            
            \item 设 $y = {\log }_{2}u,u = {2x} + 1$，则 ${y}^{\prime }_{x} = {y}^{\prime }_{u} \cdot  {u}^{\prime }_{x} = {\left( {\log }_{2}u\right) }^{\prime }{\left( 2x + 1\right) }^{\prime } = \dfrac{1}{{u}\ln 2} \times  2 = \dfrac{2}{\left( {2{x} + 1}\right) \ln 2}$，即 ${y}^{\prime } = \dfrac{2}{\left( {{2x} + 1}\right) \ln 2}$ .
            
            \item 设 $y = {e}^{u},u = {3x} + 2$ ,
            则 ${{y}}^{\prime }_{{x}} = {{y}}^{\prime }_{{u}} \cdot  {{u}}^{\prime }_{{x}} = {\left( {{e}}^{{u}}\right) }^{\prime }{\left( 3{x} + 2\right) }^{\prime } = 3{{e}}^{{u}} = 3{{e}}^{3{x} + 2}$ ,即 ${{y}}^{\prime } = 3{{e}}^{3{x} + 2}$ .
            
            \item 设 $y = \sin u,u = {2x} + \dfrac{\pi }{3}$ ,
            则 ${y}^{\prime }_x = {y}^{\prime }_u \cdot  {u}^{\prime }_x = {\left( \sin u\right) }^{\prime }{\left( 2x + \dfrac{\pi }{3}\right) }^{\prime } = 2\cos u = 2\cos \left( {{2x} + \dfrac{\pi }{3}}\right)$ ,即 ${y}^{\prime } = 2\cos \left( {{2x} + \dfrac{\pi }{3}}\right)$.

            \item ${y}^{\prime } = 2{e}^{2x+1} \sin x + {e}^{2x+1}\cos x = \left( {2\sin x + \cos x}\right) {e}^{2x+1}$ .
            
            \item ${f}^{\prime }\left( x\right)  = -\dfrac{5}{3}{\left( -x + 5\right) }^{\frac{2}{3}}$ .
            
            \item $f'\left( x\right)  = \dfrac{{2x}\left( {1 - \ln \left( {3x}\right) }\right)  + 1}{x{\left( 2x + 1\right) }^{2}}$ .
        \end{enumerate}
    \end{jieda}
    \item 已知函数 $y = f\left( x\right)$ 与 $y = g\left( x\right)$ 满足 $f\left( 1\right)  = 2,{f}^{\prime }\left( 1\right)  = 3,g\left( 1\right)  = 4$ , ${{g}}^{\prime }\left( {1}\right)  = 5$ . 对于下列函数 ${y} = {h}\left( {x}\right)$ ,求 ${h}\left( {1}\right)$ 和 ${{h}}^{\prime }\left( {1}\right)$ :
    \begin{enumerate}
        \item ${h}\left( {x}\right)  = 2{g}\left( {x}\right)  - \dfrac{1}{3}{f}\left( {x}\right)$
        \item ${h}\left( {x}\right)  = 2{g}\left( {x}\right) {f}\left( {x}\right)  - \dfrac{1}{3}$
        \item ${h}\left( {x}\right)  = \dfrac{2{g}\left( {x}\right)  - {1}}{3{f}\left( {x}\right) }$
    \end{enumerate}
    \begin{jieda}
        \begin{enumerate}
            \item $h\left( 1\right)  = \dfrac{22}{3};{h}^{\prime }\left( 1\right)  = 9$
            \item $h\left( 1\right)  = \dfrac{47}{3};{h}^{\prime }\left( 1\right)  = {44}$
            \item $h\left( 1\right)  = \dfrac{7}{6};{h}^{\prime }\left( 1\right)  =  - \dfrac{1}{12}$
        \end{enumerate}
    \end{jieda}
    \item 已知某港口一天内潮水的深度 ${y}$ (单位: ${m}$ ) 与时间 ${t}$ (单位: ${h}$ ) 近似满足函数关系 $y = 3\sin \left( {\dfrac{\pi }{12}t + \dfrac{5\pi }{6}}\right) ,0 \leq  t \leq  {24}$ . 分别求上午 6 时与下午 6 时潮水涨落的速度.
    \begin{jieda}
        ${y}^{\prime }\left( 6\right)  =  - \dfrac{\pi }{8}{m}/{h}$
        
        ${y}^{\prime }\left( {18}\right)  = \dfrac{\pi }{8}{m}/{h}$
    \end{jieda}
    
    \item 已知一罐汽水放入冰箱后的温度 $x$ (单位:℃)与时间 $t$ (单位:h)满足函数关系 $x = 4 + {16}{e}^{-{2t}}$ .
    \begin{enumerate}
        \item 求 ${x}^{\prime }\left( 1\right)$ ，并解释其实际意义
        \item 已知摄氏度 $x$ 与华氏度 $y$ (单位: $r$ )满足函数关系 $x = \dfrac{5}{9}\left( {y - {32}}\right)$ ，求 $y$ 关于 $t$ 的导数，并解释其实际意义
    \end{enumerate}
    \begin{jieda}
        \begin{enumerate}
            \item ${x}^{\prime }\left( 1\right)  =  - \dfrac{32}{{e}^{2}},t = 1$ 时汽水温度瞬时变化率,以 $\dfrac{32}{{e}^{2}}{}^{ \circ  }{C}/{h}$ 的速度下降
            
            \item ${y}^{\prime }\left( t\right)  =  - \dfrac{288}{5}{e}^{-{2t}}$ ,汽水温度华氏度在 $t$ 时刻瞬时变化率
        \end{enumerate}
    \end{jieda}
\end{enumerate}

\clearpage

\clearpage

\subsection*{导数的概念和意义(3)导数的几何意义}
\begin{enumerate}
    \item 求曲线 $y = {2}^{1 - {3x}}$ 在点 $\left( {0,2}\right)$ 处的切线方程\blkda[3]{$y = (-6ln2)x+2$}.
    
    \item 已知曲线 $y = a\ln {3x}\left( {a > 0}\right)$ 在 $x = 1$ 处的切线与两坐标轴围成的三角形的面积为 $4{\left( \ln 3 - 1\right) }^{2}$ ，则 $a =$ \blkda{8}.
    \begin{jieda}
        因为 $y = a\ln {3x}$ ,所以 $y = \dfrac{a}{x}$ ,得该曲线在 $x = 1$ 处的切线的斜率 $k = a$，切点为 $\left( {1,a\ln 3}\right)$ ,进而得切线方程为 $y - a\ln 3 = a\left( {x - 1}\right)$ .
        
        令 $y = 0$ ,得 $x = 1 - \ln 3$ ；令 $x = 0$ ,得 $y =  - a + a\ln 3$ .所以切线与两条坐标轴围成的三角形的面积 $S = \dfrac{1}{2} \times  a \times  {\left( \ln 3 - 1\right) }^{2} = 4{\left( \ln 3 - 1\right) }^{2}$，解得 $a = 8$ .
    \end{jieda}

    \item 设 $a,b \in  {R}$ ，函数 $y = \ln x + \dfrac{a}{x}$ 在 $x = 1$ 处的切线方程为 ${4y} - x - b = 0$ ， 则 $a + b =$ \blkda{$\dfrac{11}{4}$}.
    \begin{jieda}
        由 $y = \ln x + \dfrac{a}{x}$ 得 ${y}^{\prime } = \dfrac{1}{x} - \dfrac{a}{{x}^{2}}$，因为函数 $y = \ln x + \dfrac{a}{x}$ 在 $x = 1$ 处的切线方程为 ${4y} - x - b = 0$，所以 ${\left. {y}^{\prime }\right| }_{x = 1} = 1 - a = \dfrac{1}{4}$，所以 $a = \dfrac{3}{4}$

        故$y|_{x=1} = \dfrac{3}{4}$，代入切线方程，故$b = 2$，因此$a+b = \dfrac{11}{4}$
    \end{jieda}

    \item 函数 $f\left( x\right)  = \dfrac{\sqrt{3}}{3}{x}^{3} - 1$ 在 $x = 1$ 处的切线倾斜角是\blkda{$\dfrac{\pi}{3}$}.
    \begin{jieda}
        依题意, ${f}^{\prime }\left( x\right)  = \sqrt{3}{x}^{2}$ ,则函数 $f\left( x\right)  = \dfrac{\sqrt{3}}{3}{x}^{3} - 1$ 在 $x = 1$ 处的切线斜率 ${f}^{\prime }\left( 1\right)  = \sqrt{3}$，所以所求切线倾斜角为 $\dfrac{\pi }{3}$ 
    \end{jieda}

    \item 已知函数 $f\left( x\right)  = \dfrac{{{e}}^{{2x} - 2}}{x}$ ，则该函数图像在点 $\left( {1,f\left( 1\right) }\right)$ 处的切线的倾斜角的大小是\blkda{$\dfrac{\pi}{4}$}.
    \begin{jieda}
        因为 $f\left( x\right)  = \dfrac{{{e}}^{{2x} - 2}}{x}$ ,所以 ${f}^{\prime }\left( x\right)  = \dfrac{2{{e}}^{{2x} - 2}x - {{e}}^{{2x} - 2}}{{x}^{2}}$，则 ${f}^{\prime }\left( 1\right)  = \dfrac{2{{e}}^{0} - {{e}}^{0}}{1} = 1$，故倾斜角为$\dfrac{\pi}{4}$
    \end{jieda}

    \item 已知曲线 $f\left( x\right)  = 2{x}^{3} - {3x}$ ，过点 $\left( {0,0}\right)$ 作曲线的切线，则切线方程\blkda{$y = -3x$}.
    \begin{jieda}
        设切点坐标为 $\left( {{x}_{0},2{x}_{0}^{3} - 3{x}_{0}}\right)$，由 $f\left( x\right)  = 2{x}^{3} - {3x}$ ,得 ${f}^{\prime }\left( {x}_{0}\right)  = 6{x}_{0}^{2} - 3$，所以曲线 $f\left( x\right)$ 在点 $\left( {{x}_{0},2{x}_{0}^{3} - 3{x}_{0}}\right)$ 处的切线方程为 $y - \left( {2{x}_{0}^{3} - 3{x}_{0}}\right)  = \left( {6{x}_{0}^{2} - 3}\right) \left( {x - {x}_{0}}\right)$ .
        
        因为切线过点(0,0),所以 $- 2{x}_{0}^{3} + 3{x}_{0} = \left( {6{x}_{0}^{2} - 3}\right) \left( {-{x}_{0}}\right)$ ,解得 ${x}_{0} = 0$，所以切线方程为 $y =  - {3x}$ .
    \end{jieda}

    \item 已知函数 $y = f\left( x\right)$ 的导函数为 $y = {f}^{\prime }\left( x\right)$ ,且 $f\left( x\right)  = \dfrac{1}{3}{x}^{3} + {f}^{\prime }\left( 1\right) {x}^{2} + 1$ , 则 $y = f\left( x\right)$ 的图象在 $x = 3$ 处的切线方程为\blkda[3]{$y = {3x} - 8$}.
    \begin{jieda}
        函数 $f\left( x\right)  = \dfrac{1}{3}{x}^{3} + {f}^{\prime }\left( 1\right) {x}^{2} + 1$ ,求导得: ${f}^{\prime }\left( x\right)  = {x}^{2} + 2{f}^{\prime }\left( 1\right) x$ ,则 ${f}^{\prime }\left( 1\right)  = 1 + 2{f}^{\prime }\left( 1\right)$，解得 ${f}^{\prime }\left( 1\right)  =  - 1$ 
        
        因此 $f\left( x\right)  = \dfrac{1}{3}{x}^{3} - {x}^{2} + 1,{f}^{\prime }\left( x\right)  = {x}^{2} - {2x}$ ,则 $f\left( 3\right)  = 1,{f}^{\prime }\left( 3\right)  = 3$，所以所求切线方程为 $y - 1 = 3\left( {x - 3}\right)$ ,即 $y = {3x} - 8$ .
    \end{jieda}

    \item 已知曲线 $y = a{{e}}^{2x}$ 在 $x = 0$ 处的切线与直线 $x + y - 1 = 0$ 垂直，则 $a =$ \blkda{$\dfrac{1}{2}$}.
    \begin{jieda}
        因为 $y = a{{e}}^{2x}$ ,则 ${y}^{\prime } = {2a}{{e}}^{2x}$ ,可得 ${\left. {y}^{\prime }\right| }_{x = 0} = {2a}$，即曲线 $y = a{{e}}^{2x}$ 在 $x = 0$ 处的切线斜率 ${k}_{1} = {2a}$，且直线 $x + y - 1 = 0$ 的斜率 ${k}_{2} =  - 1$，由题意得 ${{k}}_{1} \cdot  {{k}}_{2} =  - {2a} =  - 1$ ,解得 $a = \dfrac{1}{2}$ .
    \end{jieda}

    \item 曲线 $y = {{e}}^{ax} + x$ 在点 $\left( {0,1}\right)$ 处的切线与直线 $x - y + 2 = 0$ 垂直，则 $a =$ \blkda{$-2$}.
    \begin{jieda}
        因为 $y = {{e}}^{ax} + x$ ,则 ${y}^{\prime } = a{{e}}^{ax} + 1,{\left. {y}^{\prime }\right| }_{x = 0} = a + 1$，又因为直线 $x - y + 2 = 0$ 的斜率为 1，由题意可得 $a + 1 =  - 1$ ,解得 $a =  - 2$ .
    \end{jieda}

    \item 若 $f\left( x\right)  = \dfrac{{{e}}^{x}}{x + 1}$ ，则曲线 $y = f\left( x\right)$ 在点 $\left( {1,f\left( 1\right) }\right)$ 处的切线与坐标轴围成的三角形面积为\blkda{$\dfrac{{e}}{8}$}.
    \begin{jieda}
        由题意得 ${f}^{\prime }\left( x\right)  = \dfrac{{{e}}^{x}\left( {x + 1}\right)  - {{e}}^{x} \cdot  1}{{\left( x + 1\right) }^{2}} = \dfrac{{{e}}^{x}\left( {x + 1 - 1}\right) }{{\left( x + 1\right) }^{2}} = \dfrac{x{{e}}^{x}}{{\left( x + 1\right) }^{2}}$，故 ${f}^{\prime }\left( 1\right)  = \dfrac{{e}}{4}$ ,又 $f\left( 1\right)  = \dfrac{{e}}{2}$，故曲线 $y = f\left( x\right)$ 在点 $\left( {1,f\left( 1\right) }\right)$ 处的切线方程为 $y - \dfrac{{e}}{2} = \dfrac{{e}}{4}\left( {x - 1}\right)$ ,即 $y = \dfrac{{e}}{4}x + \dfrac{{e}}{4}$
        
        令 $x = 0$ ,则 $y = \dfrac{{e}}{4}$ ; 令 $y = 0$ ,则 $x =  - 1$，故切线与坐标轴围成的三角形面积为 $S = \dfrac{1}{2} \times  \left| {-1}\right|  \times  \dfrac{{e}}{4} = \dfrac{{e}}{8}$
    \end{jieda}

    \item 若曲线 $y = f\left( x\right)  = 2\sin x + {2024}$ 在点 $\left( {\dfrac{\pi }{3},f\left( \dfrac{\pi }{3}\right) }\right)$ 处的切线与直线 $y = {ax} + {2024}$ 垂直，则实数 $a =$ \blkda{$-1$}.
    \begin{jieda}
        函数 $f\left( x\right)  = 2\sin x + {2024}$ ,求导得 ${f}^{\prime }\left( x\right)  = 2\cos x$，因此曲线在点 $\left( {\dfrac{\pi }{3},f\left( \dfrac{\pi }{3}\right) }\right)$ 处的切线斜率为 $k = {f}^{\prime }\left( \dfrac{\pi }{3}\right)  = 1$，而切线与直线 $y = {ax} + {2024}$ 垂直,所以 $a =  - \dfrac{1}{k} =  - 1$ .
    \end{jieda}

    \item 若直线 $y = {kx} + b$ 既与曲线 $y = {lnx} + 2$ 相切，又与曲线 $y = {{ln}\left( {x + 1}\right) }$ 相切， 则 $b =$ \blkda{$1 - \ln 2$}.
    \begin{jieda}
        设 $y = {kx} + b$ 与 $y = \ln x + 2$ 和 $y = \ln \left( {x + 1}\right)$ 的切点分别为 $\left( {{x}_{1},k{x}_{1} + b}\right) ,\left( {{x}_{2},k{x}_{2} + b}\right)$，由导数的几何意义可得 $k = \dfrac{1}{{x}_{1}} = \dfrac{1}{{x}_{2} + 1}$ ,得 ${x}_{1} = {x}_{2} + 1$，再由切点也在各自的曲线上,可得 $\left\{  \begin{array}{l} k{x}_{1} + b = \ln {x}_{1} + 2 \\  k{x}_{2} + b = \ln \left( {{x}_{2} + 1}\right)  \end{array}\right.$ ,联立上述式子解得 $\left\{  \begin{array}{l} {x}_{1} = \dfrac{1}{2} \\  k = 2 \end{array}\right.$，从而 $k{x}_{1} + b = \ln {x}_{1} + 2$ 得出 $b = 1 - \ln 2$.
    \end{jieda}

    \item 一条直线与函数 $y = \ln x$ 和 $y = {{e}}^{x}$ 的图象分别相切于点 $P\left( {{x}_{1},{y}_{1}}\right)$ 和点 $Q\left( {{x}_{2},{y}_{2}}\right)$ ，则 $\left( {1 - {x}_{1}}\right) \left( {1 + {x}_{2}}\right)$ 的值为\blkda{2}.
    \begin{jieda}
        对 $y = \ln x$ 求导得: ${y}^{\prime } = \dfrac{1}{x}$，则在 $P\left( {{x}_{1},{y}_{1}}\right)$ 处切线斜率为 $\dfrac{1}{{x}_{1}}$ ,且 ${y}_{1} = \ln {x}_{1}$，对于 $y = {{e}}^{x}$ 求导得: ${y}^{\prime } = {{e}}^{x}$，则在 $Q\left( {{x}_{2},{y}_{2}}\right)$ 处切线斜率为 ${{e}}^{{x}_{2}}$ ,且 ${y}_{2} = {{e}}^{{x}_{2}}$
        
        由题意可得: $\dfrac{1}{{x}_{1}} = {{e}}^{{x}_{2}}$ ,即 ${x}_{2} = \ln \dfrac{1}{{x}_{1}}$，又切线斜率 $k = \dfrac{{y}_{2} - {y}_{1}}{{x}_{2} - {x}_{1}} = \dfrac{{{e}}^{{x}_{2}} - \ln {x}_{1}}{{x}_{2} - {x}_{1}} = \dfrac{\dfrac{1}{{x}_{1}} - \ln {x}_{1}}{\ln \dfrac{1}{{x}_{1}} - {x}_{1}} = \dfrac{1}{{x}_{1}}$，可得: $1 - {x}_{1}\ln {x}_{1} =  - \ln {x}_{1} - {x}_{1}$ ,即 $\left( {{x}_{1} - 1}\right) \ln {x}_{1} = 1 + {x}_{1}$，所以 $\left( {1 - {x}_{1}}\right) \left( {1 + {x}_{2}}\right)  = \left( {1 - {x}_{1}}\right) \left( {1 - \ln {x}_{1}}\right)  = 1 - {x}_{1} + \left( {{x}_{1} - 1}\right) \ln {x}_{1} = 1 - {x}_{1} + 1 + {x}_{1} = 2$
    \end{jieda}

    \item 直线 $l : y = {kx} + b$ 是曲线 $y = \ln x$ 和 $y = {{e}}^{x - 2}$ 的公切线，则 $k + b =$ \blkda{$\dfrac{1}{e}$或0}.
    \begin{jieda}
        对于 $y = \ln x$ ,设切点为 $\left( {{x}_{1},\ln {x}_{1}}\right)$ ,求导得 ${y}^{\prime } = \dfrac{1}{x}$，则在该点处的斜率为 $k = \dfrac{1}{{x}_{1}}$，则切线方程为: $y - \ln {x}_{1} = \dfrac{1}{{x}_{1}}\left( {x - {x}_{1}}\right)$ ,即 $y = \dfrac{1}{{x}_{1}}x + \ln {x}_{1} - 1$
        
        对于 $y = {{e}}^{x - 2}$ ,设切点为 $\left( {{x}_{2},{{e}}^{{x}_{2} - 2}}\right)$ ,求导得 ${y}^{\prime } = {{e}}^{x - 2}$，则在该点处的斜率为 $k = {{e}}^{{x}_{2} - 2}$，则切线方程为: $y - {{e}}^{{x}_{2} - 2} = {{e}}^{{x}_{2} - 2}\left( {x - {x}_{2}}\right)$ ,即 $y = {{e}}^{{x}_{2} - 2}x + \left( {1 - {x}_{2}}\right) {{e}}^{{x}_{2} - 2}$
        
        因为 $l : y = {kx} + b$ 是公切线，所以 $\left\{  \begin{array}{l} \dfrac{1}{{x}_{1}} = {{e}}^{{x}_{2} - 2} \\  \ln {x}_{1} - 1 = \left( {1 - {x}_{2}}\right) {{e}}^{{x}_{2} - 2} \end{array}\right.$ ,即 ${x}_{1} = {{e}}^{2 - {x}_{2}}$，所以 $\ln {{e}}^{2 - {x}_{2}} - 1 = \left( {1 - {x}_{2}}\right) {{e}}^{{x}_{2} - 2}$ ,即 $1 - {x}_{2} = \left( {1 - {x}_{2}}\right) {{e}}^{{x}_{2} - 2}$，所以 $\left( {1 - {x}_{2}}\right) \left( {{{e}}^{{x}_{2} - 2} - 1}\right)  = 0$，即 $1 - {x}_{2} = 0$ 或 ${x}_{2} - 2 = 0$ ,解得 ${x}_{2} = 1$ 或 ${x}_{2} = 2$
        
        当 ${x}_{2} = 1$ 时,此时 $k = {{e}}^{{x}_{2} - 2} = \dfrac{1}{{e}},b = \left( {1 - {x}_{2}}\right) {{e}}^{{x}_{2} - 2} = 0$ ,所以 $k + b = \dfrac{1}{{e}}$
        
        当 ${x}_{2} = 2$ 时,此时 $k = {{e}}^{{x}_{2} - 2} = 1,b = \left( {1 - {x}_{2}}\right) {{e}}^{{x}_{2} - 2} =  - 1$ ,所以 $k + b = 0$
        
        所以 $k + b = 0$ 或 $\dfrac{1}{{e}}$
    \end{jieda}

    \item 已知函数 $f\left( x\right)  = \ln x - \dfrac{1}{x}$ ，直线 $y = {mx} + n$ 是曲线 $y = f\left( x\right)$ 的一条切线， 则 $m + {2n}$ 的取值范围是\blkda[3]{$\lbrack  - 2\ln 2 - 4, + \infty )$}.
    \begin{jieda}
        由 $f\left( x\right)  = \ln x - \dfrac{1}{x},x > 0$ 可得 ${f}^{\prime }\left( x\right)  = \dfrac{1}{x} + \dfrac{1}{{x}^{2}}$，设切点为 $P\left( {t,f\left( t\right) }\right)$ ,则 $k = {f}^{\prime }\left( t\right)  = \dfrac{1}{t} + \dfrac{1}{{t}^{2}}$，所以曲线 $y = f\left( x\right)$ 在切点 $P\left( {t,f\left( t\right) }\right)$ 处的切线方程为 $y - f\left( t\right)  = {f}^{\prime }\left( t\right) \left( {x - t}\right)$，整理得 $y = \left( {\dfrac{1}{t} + \dfrac{1}{{t}^{2}}}\right) x + \ln t - \dfrac{2}{t} - 1$ 
        
        所以 $m + {2n} = \dfrac{1}{t} + \dfrac{1}{{t}^{2}} + 2\left( {\ln t - \dfrac{2}{t} - 1}\right)  = \dfrac{1}{{t}^{2}} + 2\ln t - \dfrac{3}{t} - 2$ 
        
        令 $g\left( x\right)  = \dfrac{1}{{x}^{2}} + 2\ln x - \dfrac{3}{x} - 2\left( {x > 0}\right)$ ,则 ${g}^{\prime }\left( x\right)  =  - \dfrac{2}{{x}^{3}} + \dfrac{2}{x} + \dfrac{3}{{x}^{2}} = \dfrac{2{x}^{2} + {3x} - 2}{{x}^{3}}$ 
        
        令 ${g}^{\prime }\left( x\right)  = 0$ ,解得 $x = \dfrac{1}{2}$ 
        
        当 $0 < x < \dfrac{1}{2}$ 时, ${g}^{\prime }\left( x\right)  < 0,g\left( x\right)$ 单调递减，当 $x > \dfrac{1}{2}$ 时, ${g}^{\prime }\left( x\right)  > 0,g\left( x\right)$ 单调递增
        
        故 $g{\left( x\right) }_{\min } = g\left( \dfrac{1}{2}\right)  =  - 2\ln 2 - 4$，则 $m + {2n}$ 的取值范围是 $\lbrack  - 2\ln 2 - 4, + \infty )$
    \end{jieda}

    \item 
    \begin{enumerate}
        \item 设 $a$ 、 $b$ 为实数，比较 ${a}^{2} + {b}^{2}$ 与 ${2a} - {2b} - 2$ 的值的大小
        \item 已知 $f\left( x\right)  = 3{x}^{2}$ ，求曲线 $y = f\left( x\right)$ 在点 $P\left( {-1,3}\right)$ 处的切线方程.
    \end{enumerate}
    \begin{jieda}
        \begin{enumerate}
            \item  ${a}^{2} + {b}^{2} - \left( {{2a} - {2b} - 2}\right)  = \left( {{a}^{2} - {2a} + 1}\right)  + \left( {{b}^{2} + {2b} + 1}\right)  = {\left( a - 1\right) }^{2} + {\left( b + 1\right) }^{2} \geq  0$，当且仅当 $a = 1,b =  - 1$ 时等号成立，所以 ${a}^{2} + {b}^{2} \geq  {2a} - {2b} - 2$ 
            
            \item ${f}^{\prime }\left( x\right)  = {6x}$ ,则 ${f}^{\prime }\left( {-1}\right)  =  - 6$，因此,曲线 $y = 3{x}^{2}$ 在点 $P\left( {-1,3}\right)$ 处的切线斜率为 $-6$，于是,所求切线方程为 $y - 3 =  - 6\left( {x + 1}\right)$ ,即 ${6x} + y + 3 = 0$ .
        \end{enumerate}
    \end{jieda}
\end{enumerate}

\clearpage

\subsection*{导数的应用(1)——函数的单调性}
A组：
\begin{enumerate}
    \item 函数 $f\left( x\right)  = {x}^{2} - {x}^{3}$ 的严格递增区间为\blkda{$\left(0, \dfrac{2}{3}\right)$}，严格递减区间为\blkda[3]{$(-\infty, 0), \left(\dfrac{2}{3}, +\infty \right)$}.

    \item 函数 $f\left( x\right)  = {e}^{x} - x$ 的单调减区间是\blkda{$(-\infty, 0)$}.

    \item 函数 $y = 4{x}^{2} + \dfrac{1}{x}$ 的单调递增区间是\blkda{$\left(\dfrac{1}{2}, +\infty\right)$}.

    \item 已知函数 $f\left( x\right)  = \dfrac{1}{3}{x}^{3} + b{x}^{2} + \left( {b + 2}\right) x + 3$ 在 $\mathbf{R}$ 上单调递增,则实数 $b$ 的取值范围是\blkda{$[-1, 2]$}.
    \begin{jieda}
        $f'(x) = x^2+2bx+(b+2) \geq 0$恒成立，$\Delta = 4b^2 - 4(b+2) < 0$，解得$b \in [-1, 2]$
    \end{jieda}

    \item 证明: 函数 $y = {x}^{3} + {4x}$ 在 $\left( {-\infty , + \infty }\right)$ 上严格增.
    \begin{jieda}
        $y' = 3x^2 + 4$，故$y' > 0$恒成立.
    \end{jieda}

    \item 求函数 $f\left( x\right)  = 2{x}^{3} - 6{x}^{2} - {18x} - 7$ 的单调区间.
    \begin{jieda}
        $f'(x) = 6x^2-12x-18 = 6(x^2-2x-3)$

        单调增区间为 $\left( {-\infty , - 1}\right)$ , $\left( {3, + \infty }\right)$ ,单调减区间为$(-1,3)$
    \end{jieda}
\end{enumerate}

B组：
\begin{enumerate}
    \item 已知 $f\left( x\right)  = {x}^{3} - {ax}$ 在 $\left( {1, + \infty }\right)$ 上是严格增函数，则 $a$ 的最大值是\blkda{3}.
    \begin{jieda}
        $f'(x) = 3x^2-a \geq 0$在$\left( {1, + \infty }\right)$ 上恒成立.
    \end{jieda}

    \item 函数 $y = \dfrac{1}{2}{x}^{2} - \ln x$ 的单调减区间为\blkda{$(0, 1)$}.

    \item 设 $f\left( x\right)  = 4{x}^{3} + m{x}^{2} + \left( {m - 3}\right) x + n\left( {m,n \in  {R}}\right)$ 是 ${R}$ 上的单调增函数,则 $m =$ \blkda{6}.
    \begin{jieda}
        由题意可知, ${f}^{\prime }\left( x\right)  = {12}{x}^{2} + {2mx} + m - 3$
        
        因为 $f\left( x\right)$ 是 ${R}$ 上的单调增函数，所以 ${f}^{\prime }\left( x\right)  = {12}{x}^{2} + {2mx} + m - 3 \geq  0$ 在 $\mathbf{R}$ 上恒成立，所以 ${\left( 2m\right) }^{2} - 4 \times  {12} \times  \left( {m - 3}\right)  \leq  0$ ,即 ${\left( m - 6\right) }^{2} \leq  0$ ,解得 $m = 6$
    \end{jieda}

    \item 函数 $f\left( x\right)  = \left( {x - 3}\right) {{e}}^{x}$ 的单调增区间是\blkda{$(2, +\infty)$}.

    \item 请根据图中的函数图像，将下列数值按从小到大的顺序排列:\blkda{\ding{174}\ding{176}\ding{173}\ding{175}\ding{177}\ding{172}}.
    \begin{center}
    \includegraphics[width=0.2\textwidth]{bodyxb/src/bo_d6jmltk601uc73e0s4bg_0_1065_1698_340_213_0.jpg}
    \end{center}
    \begin{dingautolist}{172}
        \item 曲线在 A 处切线的斜率
        \item 曲线在 B 处切线的斜率
        \item 曲线在 C 处切线的斜率
        \item 割线 AB 的斜率
        \item 数值 0
        \item 数值 1
    \end{dingautolist}
    
    \item 设函数 $f\left( x\right)  = {ax} - \dfrac{a}{x} - 2\ln x$ 的导函数为 ${f}^{\prime }\left( x\right)$ ,且 ${f}^{\prime }\left( 2\right)  = 0$ ,则 $f\left( x\right)$ 的单调递减区间为\blkda{$\left( {\dfrac{1}{2},2}\right)$}.
    \begin{jieda}
        $f\left( x\right)$ 定义域为 $\left( {0, + \infty }\right) ,{f}^{\prime }\left( x\right)  = a + \dfrac{a}{{x}^{2}} - \dfrac{2}{x}$，所以 ${f}^{\prime }\left( 2\right)  = a + \dfrac{a}{4} - 1 = 0$，解得 $a = \dfrac{4}{5}$，所以 $f\left( x\right)  = \dfrac{4}{5}x - \dfrac{4}{5x} - 2\ln x,{f}^{\prime }\left( x\right)  = \dfrac{4}{5} + \dfrac{4}{5{x}^{2}} - \dfrac{2}{x} = \dfrac{4{x}^{2} - {10x} + 4}{5{x}^{2}}$，由 ${f}^{\prime }\left( x\right)  < 0$ 解得 $\dfrac{1}{2} < x < 2$，所以 $f\left( x\right)$ 的单调递减区间为 $\left( {\dfrac{1}{2},2}\right)$
    \end{jieda}

    \item 若函数 $f\left( x\right)  = {x}^{3} - {3a}{x}^{2} - 9{a}^{2}x + 1$ 在 $\left( {-1,2}\right)$ 单调递减，则 $a$ 的取值范围是\blkda[3]{ $\left( {-\infty , - 2\rbrack \cup \lbrack 1, + \infty }\right)$}.
    \begin{jieda}
        ${f}^{\prime }\left( x\right)  = 3{x}^{2} - {6ax} - 9{a}^{2}$
        
        令 ${f}^{\prime }\left( x\right)  = 0$ ,解得 $x = {3a}$ 或 $x =  - a$
        
        若函数 $f\left( x\right)  = {x}^{3} - {3a}{x}^{2} - 9{a}^{2}x + 1$ 在$(-1,2)$单调递减,则 ${f}^{\prime }\left( x\right)  \leq  0$ 在$(-1,2)$恒成立

        易知$a = 0$不合题意.
        
        当 $a < 0$ 时, ${3a} <  - a,\therefore \left\{  {\begin{array}{l}  - 1 \geq  {3a} \\   - a \geq  2 \end{array} \Longrightarrow  \left\{  {\begin{array}{l} a \leq   - \dfrac{1}{3} \\  a \leq   - 2 \end{array} \Longrightarrow  a \leq   - 2}\right. }\right.$
        
        当 $a > 0$ 时, ${3a} >  - a,\therefore \left\{  {\begin{array}{l}  - a \leq   - 1 \\  {3a} \geq  2 \end{array} \Longrightarrow  \left\{  {\begin{array}{l} a \geq  1 \\  a \geq  \dfrac{2}{3} \end{array} \Longrightarrow  a \geq  1}\right. }\right.$ ,
        
        故$a$ 的取值范围为 $\left( {-\infty , - 2\rbrack \cup \lbrack 1, + \infty }\right)$ .
    \end{jieda}

    \item 若函数 $f\left( x\right)  = \cos {2x} + a\sin x$ 在区间 $\left( {\dfrac{\pi }{6},\dfrac{\pi }{2}}\right)$ 上是减函数，则 $a$ 的取值范围是\blkda{$(-\infty, 2]$}.
    \begin{jieda}
        因为 ${f}^{\prime }\left( x\right)  =  - 2\sin {2x} + a\cos x =  - 4\sin x\cos x + a\cos x = \cos x\left( {-4\sin x + a}\right)$，而 $x \in  \left( {\dfrac{\pi }{6},\dfrac{\pi }{2}}\right)$ 时,函数 $f\left( x\right)$ 单调递减,所以 ${f}^{\prime }\left( x\right)  < 0$ 在 $x \in  \left( {\dfrac{\pi }{6},\dfrac{\pi }{2}}\right)$ 恒成立，即 $\cos x\left( {-4\sin x + a}\right)  < 0$ 恒成立
        
        因为 $\cos x > 0$ ,所以 $- 4\sin x + a < 0$，即 $a < 4\sin x$ 在 $x \in  \left( {\dfrac{\pi }{6},\dfrac{\pi }{2}}\right)$ 恒成立，因为 $y = 4\sin x$ 在 $\left( {\dfrac{\pi }{6},\dfrac{\pi }{2}}\right)$ 上单调递增，则 $4\sin x > 4\sin \dfrac{\pi }{6} = 2$ ,所以 $a \leq  2$ .
    \end{jieda}

    \item 若函数 $f\left( x\right)  = \dfrac{2}{3}{x}^{3} - 2{x}^{2} + \left( {4 - {2a}}\right) x + 3$ 在区间 $\left\lbrack  {0,3}\right\rbrack$ 内存在单调递增区间，则实数 $a$ 的取值范围是\blkda{$\left( {-\infty ,5}\right)$}.
    \begin{jieda}
        求导得: ${f}^{\prime }\left( x\right)  = 2{x}^{2} - {4x} + \left( {4 - {2a}}\right)$，$f\left( x\right)$ 区间 $\left\lbrack  {0,3}\right\rbrack$ 内存在单调递增区间,也就是 $2{x}^{2} - {4x} + 4 - {2a} \geq 0$ 在 $\left\lbrack  {0,3}\right\rbrack$ 内有解区间,整理得$a \leq {x}^{2} - {2x} + 2$ 在 $\left\lbrack  {0,3}\right\rbrack$ 内有解区间
                
        令 $g\left( x\right)  = {x}^{2} - {2x} + 2$ ,其对称轴为 $x = 1$ ,在区间 $\left\lbrack  {0,3}\right\rbrack$ 内,
        计算端点值: $g\left( 0\right)  = {0}^{2} - 2 \times  0 + 2 = 2,g\left( 3\right)  = {3}^{2} - 2 \times  3 + 2 = 5$ ,
        所以 $g\left( x\right)$ 在 $\left\lbrack  {0,3}\right\rbrack$ 上的最大值为 5
        
        因此实数 $a$ 的取值范围是 $\left( {-\infty ,5}\right)$ .
    \end{jieda}

    \item 已知函数 $f\left( x\right)  = a\ln x - {x}^{2}$ 存在单调递增区间,则实数 $a$ 的取值范围为\blkda{$(0, +\infty)$}.
    \begin{jieda}
        易知 $f\left( x\right)$ 的定义域为 $\left( {0, + \infty }\right)$ ,又 ${f}^{\prime }\left( x\right)  = \dfrac{a}{x} - {2x} = \dfrac{a - 2{x}^{2}}{x}$ 
        
        由题意可知 $a - 2{x}^{2} > 0$ 在 $\left( {0, + \infty }\right)$ 上有解,即 $a > 2{x}^{2}$ 在 $\left( {0, + \infty }\right)$ 上有解，可得 $a > {\left( 2{x}^{2}\right) }_{\min },x \in  \left( {0, + \infty }\right)$ ,所以 $a > 0$ .
    \end{jieda}

    \item 已知函数 $f\left( x\right)  = {{e}}^{ax}\sin x$ . 当 $a = 1$ 时，求 $f\left( x\right)$ 的单调递增区间\blkda[5]{$\left\lbrack  {-\dfrac{\pi }{4} + {2k\pi },\dfrac{3\pi }{4} + {2k\pi }}\right\rbrack  ,k \in  \mathbf{Z}$}.
    \begin{jieda}
        当 $a = 1$ 时, $f\left( x\right)  = {{e}}^{x}\sin x$ ,则 ${f}^{\prime }\left( x\right)  = {{e}}^{x}\left( {\sin x + \cos x}\right)  = \sqrt{2}{{e}}^{x}\sin \left( {x + \dfrac{\pi }{4}}\right)$ 
        
        令 ${f}^{\prime }\left( x\right)  \geq  0$ ,得到 ${2k\pi } \leq  x + \dfrac{\pi }{4} \leq  \pi  + {2k\pi },k \in  Z$，解得: $x \in  \left\lbrack  {-\dfrac{\pi }{4} + {2k\pi },\dfrac{3\pi }{4} + {2k\pi }}\right\rbrack  ,k \in  \mathbf{Z}$
    \end{jieda}
    
    \item 已知函数 $y = f\left( x\right)$ 的导函数 $y = {f}^{\prime }\left( x\right)$ 的图象如图所示,则该函数的图象可能是\dotfill{(\kh{B})}
    
    \begin{center}
    \includegraphics[width=0.2\textwidth]{bodyxb/src/bo_d6jmltk601uc73e0s4bg_1_235_1130_268_279_0.jpg}
    \end{center}

    \xx{\includegraphics[width=0.35\textwidth]{bodyxb/src/bo_d6jmltk601uc73e0s4bg_1_607_1255_274_263_0.jpg}}{\includegraphics[width=0.35\textwidth]{bodyxb/src/bo_d6jmltk601uc73e0s4bg_1_1000_1253_267_263_0.jpg}}{\includegraphics[width=0.35\textwidth]{bodyxb/src/bo_d6jmltk601uc73e0s4bg_1_599_1518_274_269_0.jpg}}{\includegraphics[width=0.35\textwidth]{bodyxb/src/bo_d6jmltk601uc73e0s4bg_1_1005_1515_267_270_0.jpg}}


    \item 已知 $p : \dfrac{{3m} - 1}{m - 1} \leq  0,q$ : 函数 $f\left( x\right)  = \dfrac{1}{3}{x}^{3} - \dfrac{3m}{2}{x}^{2} + 1$ 在区间 $\left( {1,3}\right)$ 上不单调, 则 $p$ 是 $q$ 的\dotfill{(\kh{B})}
    \xx{充分不必要条件}{必要不充分条件}{充要条件}{既不充分也不必要条件}
    \begin{jieda}
        由 $\dfrac{{3m} - 1}{m - 1} \leq  0$ ,可化为 $\left\{  \begin{array}{l} \left( {{3m} - 1}\right) \left( {m - 1}\right)  \leq  0 \\  m - 1 \neq  0 \end{array}\right.$，得 $\dfrac{1}{3} \leq  m < 1$
        
        ${f}^{\prime }\left( x\right)  = {x}^{2} - {3mx} = x\left( {x - {3m}}\right)$，要使函数 $f\left( x\right)$ 在区间(1,3)上不单调. 则有 $1 < {3m} < 3$ ,解得 $\dfrac{1}{3} < m < 1$
        
        所以 $p$ 是 $q$ 的必要不充分条件.
    \end{jieda}

    \item 已知 $f\left( x\right)  = a\left( {x - \ln x}\right)  + \dfrac{{2x} - 1}{{x}^{2}},a > 0$ . 讨论 $f\left( x\right)$ 的单调性.
    \begin{jieda}
        $f'(x) = a\left(1-\dfrac{1}{x}\right) - \dfrac{2}{x^2}+\dfrac{2}{x^3} = a\left(1-\dfrac{1}{x}\right)-\dfrac{2}{x^2}\left(1-\dfrac{1}{x}\right) = \left(a-\dfrac{2}{x^2}\right)\left(1-\dfrac{1}{x}\right), x > 0$
        \begin{dingautolist}{172}
            \item 当 $0 < a < 2$ 时,单调增区间为 $\left( {0,1}\right) ,\left( {\sqrt{\dfrac{2}{a}}, + \infty }\right)$ ,单调减区间为 $\left( {1,\sqrt{\dfrac{2}{a}}}\right)$
            \item 当 $a = 2$ 时,单调增区间为 $\left( {0, + \infty }\right)$
            \item 当 $a > 2$ 时,单调增区间为 $\left( {0,\sqrt{\dfrac{2}{a}}}\right) ,\left( {1, + \infty }\right)$ ,单调减区间为 $\left( {\sqrt{\dfrac{2}{a}},1}\right)$
        \end{dingautolist}
    \end{jieda}

    \item 已知函数 $f\left( x\right)  = {x}^{3} + a{x}^{2} + 1,a \in  {R}$ . 若函数 $f\left( x\right)$ 在区间 $\left( {-\dfrac{2}{3},0}\right)$ 内严格减,求 $a$ 的取值范围.
    \begin{jieda}
        ${f}^{\prime }\left( x\right)  = 3{x}^{2} + {2ax} = {3x}\left( {x + \dfrac{2a}{3}}\right) .$
        \begin{dingautolist}{172}
            \item 当 $a = 0$ 时, ${f}^{\prime }\left( x\right)  = 3{x}^{2} \geq  0$ 恒成立,
            $\therefore f\left( x\right)  = {x}^{3} + 1$ 在 $\mathbf{R}$ 上严格增
            \item 当 $a > 0$ 时,令 ${f}^{\prime }\left( x\right)  > 0$ ,得 $x \in  \left( {-\infty , - \dfrac{2a}{3}}\right)  \cup  \left( {0, + \infty }\right)$ ；
            令 ${f}^{\prime }\left( x\right)  < 0$ ,得 $x \in  \left( {-\dfrac{2a}{3},0}\right)$ ,
            所以 $f\left( x\right)$ 的单调增区间为 $\left( {-\infty , - \dfrac{2a}{3}}\right)$ 及 $\left( {0, + \infty }\right)$ ,单调减区间为 $\left( {-\dfrac{2a}{3},0}\right)$
            \item 当 $a < 0$ 时,令 ${f}^{\prime }\left( x\right)  > 0$ ,得 $x \in  \left( {-\infty ,0}\right)  \cup  \left( {-\dfrac{2a}{3}, + \infty }\right)$ ；
            令 ${f}^{\prime }\left( x\right)  < 0$ ,得 $x \in  \left( {0, - \dfrac{2a}{3}}\right)$ ,
            所以 $f\left( x\right)$ 的单调增区间为 $\left( {-\infty ,0}\right)$ 及 $\left( {-\dfrac{2a}{3}, + \infty }\right)$ ,单调减区间为 $\left( {0, - \dfrac{2a}{3}}\right)$ 
        \end{dingautolist}
        综上,当 $a > 0$ 时, $\left( {-\dfrac{2}{3},0}\right)  \subseteq  \left( {-\dfrac{2a}{3},0}\right)$ ,所以 $- \dfrac{2a}{3} \leq   - \dfrac{2}{3}$ ,解得 $a \geq  1$ ,即 $a$ 的取值范围为 $\lbrack 1, + \infty )$ .
    \end{jieda}

    \item 已知函数 $f\left( x\right)  = \dfrac{{f}^{\prime }\left( 3\right) }{8} \cdot  {x}^{2} + {ax} - 3\ln x\left( {a \in  \mathbf{R}}\right)$ 在 $x = 1$ 处的切线与 $y$ 轴垂直.
    \begin{enumerate}
        \item 求函数 $f\left( x\right)$ 的解析式
        \item 求函数 $f\left( x\right)$ 的单调区间
    \end{enumerate}
    \begin{jieda}
        \begin{enumerate}
            \item 由题意得 ${f}^{\prime }\left( x\right)  = \dfrac{{f}^{\prime }\left( 3\right) }{4} \cdot  x + a - \dfrac{3}{x}$ ,令 $x = 3$ ,得 ${f}^{\prime }\left( 3\right)  = \dfrac{{f}^{\prime }\left( 3\right) }{4} \cdot  3 + a - \dfrac{3}{3}$ ,解得
            ${f}^{\prime }\left( 3\right)  = {4a} - 4,$
            又函数在 $x = 1$ 处的切线与 $y$ 轴垂直, $\therefore {f}^{\prime }\left( 1\right)  = 0$ ,
            $\therefore {f}^{\prime }\left( 1\right)  = \dfrac{4\left( {a - 1}\right) }{4} \cdot  1 + a - 3 = {2a} - 4 = 0,\therefore a = 2$ ,则 ${f}^{\prime }\left( 3\right)  = 4$ ,
            $\therefore$ 函数 $f\left( x\right)$ 的解析式为 $f\left( x\right)  = \dfrac{1}{2}{x}^{2} + {2x} - 3\ln x$ .
            \item 由( 1 )可知 $f\left( x\right)  = \dfrac{1}{2} \cdot  {x}^{2} + {2x} - 3\ln x$ ,
            则 ${f}^{\prime }\left( x\right)  = x + 2 - \dfrac{3}{x} = \dfrac{{x}^{2} + {2x} - 3}{x} = \dfrac{\left( {x + 3}\right) \left( {x - 1}\right) }{x}$ ,
            又函数 $f\left( x\right)$ 的定义域为 $x > 0,\therefore x + 3 > 0$ ,
            故当 $x - 1 > 0$ 时, $x > 1$ ,此时 ${f}^{\prime }\left( x\right)  > 0$ ,函数 $f\left( x\right)$ 在 $\left( {1, + \infty }\right)$ 上单调递增;
            当 $x - 1 < 0$ 时, $0 < x < 1$ ,此时 ${f}^{\prime }\left( x\right)  < 0$ ,函数 $f\left( x\right)$ 在(0,1)上单调递减.
            所以 $f\left( x\right)$ 的单调增区间为 $\left( {1, + \infty }\right)$ ,单调减区间为(0,1).
        \end{enumerate}
    \end{jieda}
    
    \item 已知 $f\left( x\right)  = {{e}}^{x} - {ax} - 1$ .
    \begin{enumerate}
        \item 当 $a = 2$ 时，讨论 $f\left( x\right)$ 的单调区间
        \item 若 $f\left( x\right)$ 在定义域 ${R}$ 内单调递增，求 $a$ 的取值范围
    \end{enumerate}
    \begin{jieda}
        \begin{enumerate}
            \item 当 $a = 2$ 时, ${f}^{\prime }\left( x\right)  = {e}^{x} - 2$ .
            令 ${f}^{\prime }\left( x\right)  > 0$ ,即 ${e}^{x} - 2 > 0$ 解得： $x > \ln 2$ ；
            令 ${f}^{\prime }\left( x\right)  < 0$ ,即 ${e}^{x} - 2 < 0$ 解得： $x < \ln 2$ ；
            $\therefore$ 当 $a = 2$ 时,函数 $f\left( x\right)$ 的单调增区间是 $\left( {\ln 2, + \infty }\right)$ ,递减区间为 $\left( {-\infty ,\ln 2}\right)$ .
            \item $\because f\left( x\right)  = {e}^{x} - {ax} - 1,x \in  \mathbf{R}\therefore {f}^{\prime }\left( x\right)  = {e}^{x} - a$
            $\because f\left( x\right)$ 在 ${R}$ 上单调递增,即 ${f}^{\prime }\left( x\right)  = {e}^{x} - a \geq  0$ 恒成立,
            $\because x \in  {R}$ 时 ${e}^{x} \in  \left( {0, + \infty }\right)$
            $\therefore a \leq  0$ ,即 $a$ 的取值范围为 $( - \infty ,0\rbrack$ .
        \end{enumerate}
    \end{jieda}
    
    \item 设函数 $f\left( x\right)  = {ax} - 2 - \ln x\left( {a \in  {R}}\right)$ .
    \begin{enumerate}
        \item 若 $f\left( x\right)$ 在点 $\left( {{e},f\left( {e}\right) }\right)$ 处的切线为 $x - {e}y + b = 0$ ，求 $a,b$ 的值
        \item 求 ${f}^{\prime }\left( x\right)  < 0$ 的解集
    \end{enumerate}
    \begin{jieda}
        \begin{enumerate}
            \item 易知 $f\left( x\right)  = {ax} - 2 - \ln x\left( {a \in  \mathbf{R}}\right)$ 的定义域为 $\left( {0, + \infty }\right)$ ,
            因为 ${f}^{\prime }\left( x\right)  = a - \dfrac{1}{x}$ ,
            因为 $f\left( x\right)$ 在点 $\left( {{e},f\left( {e}\right) }\right)$ 处的切线为 $x - {e}y + b = 0$ ,
            所以 ${f}^{\prime }\left( {e}\right)  = a - \dfrac{1}{{e}} = \dfrac{1}{{e}}$ ,所以 $a = \dfrac{2}{{e}}$ ,所以 $f\left( {e}\right)  =  - 1$ ,
            把点$(e, -1)$代入 $x - {e}y + b = 0$ 得： $b =  - 2{e}$ .
            即 $a,b$ 的值为: $a = \dfrac{2}{{e}},b =  - 2{e}$ .
            \item ${f}^{\prime }\left( x\right)  = a - \dfrac{1}{x} = \dfrac{{ax} - 1}{x}\left( {x > 0}\right)$ .
            \begin{dingautolist}{172}
                \item 当 $a \leq  0$ 时, ${f}^{\prime }\left( x\right)  < 0$ 在 $\left( {0, + \infty }\right)$ 上恒成立,所以 ${f}^{\prime }\left( x\right)  < 0$ 的解集为 $\left( {0, + \infty }\right)$
                \item 当 $a > 0$ 时,令 ${f}^{\prime }\left( x\right)  < 0$ ,解得： $0 < x < \dfrac{1}{a}$ 
            \end{dingautolist}
            
            综上所述: 当 $a \leq  0$ 时, ${f}^{\prime }\left( x\right)  < 0$ 的解集为 $\left( {0, + \infty }\right)$ ;
            当 $a > 0$ 时, ${f}^{\prime }\left( x\right)  < 0$ 的解集为 $\left( {0,\dfrac{1}{a}}\right)$.
        \end{enumerate}
    \end{jieda}
\end{enumerate}
\clearpage
\subsection*{导数的应用(2)——极值}
A组：
\begin{enumerate}
    \item 函数 $f\left( x\right)  = \ln x - x$ 的极大值点为\blkda{1}.
    \begin{jieda}
        因为 $f\left( x\right)  = \ln x - x,\left( {x > 0}\right)$ ,所以 ${f}^{\prime }\left( x\right)  = \dfrac{1}{x} - 1 = \dfrac{1 - x}{x}$ 
        \begin{dingautolist}{172}
            \item 当 $x > 1$ 时, ${f}^{\prime }\left( x\right)  < 0$ ,函数 $f\left( x\right)$ 单调递减
            \item 当 $0 < x < 1$ 时, ${f}^{\prime }\left( x\right)  > 0$ ,函数 $f\left( x\right)$ 单调递增
        \end{dingautolist}
        所以在 $x = 1$ 处取得极大值,即函数 $f\left( x\right)  = \ln x - x$ 的极大值点为 1 .
    \end{jieda}
    
    \item 函数 $f\left( x\right)  = {{e}}^{x + 1} - x$ 的极小值点为\blkda{$-1$}.
    \begin{jieda}
        由 $f\left( x\right)  = {{e}}^{x + 1} - x$ 可得 ${f}^{\prime }\left( x\right)  = {{e}}^{x + 1} - 1$，令 ${f}^{\prime }\left( x\right)  = 0$ 可得 ${{e}}^{x + 1} - 1 = 0$ ,即 ${{e}}^{x + 1} = 1 = {{e}}^{0}$ ,解得 $x =  - 1$ 
        \begin{dingautolist}{172}
            \item 当 $x \in  \left( {-\infty , - 1}\right)$ 时, ${f}^{\prime }\left( x\right)  < 0$ ,即 $f\left( x\right)$ 单调递减
            \item 当 $x \in  \left( {-1, + \infty }\right)$ 时, ${f}^{\prime }\left( x\right)  > 0$ ,即 $f\left( x\right)$ 单调递增
        \end{dingautolist}
        所以 $x =  - 1$ 是 $f\left( x\right)$ 的极小值点.
    \end{jieda}

    \item 函数 $f\left( x\right)  = \dfrac{1}{3}{x}^{3} - {4x} + 4$ 的极大值为\blkda{$\dfrac{28}{3}$}.

    \item 已知函数 $f\left( x\right)$ 的导函数 ${f}^{\prime }\left( x\right)$ 的图象如图所示，则函数 $f\left( x\right)$ 的极值点的个数有\blkda{1}个.

    \begin{center}
    \includegraphics[width=0.2\textwidth]{bodyxb/src/bo_d6kkqnk601uc73e1cnh0_0_213_495_350_344_0.jpg}
    \end{center}
\end{enumerate}

B组：
\begin{enumerate}
    \item 若函数 $f\left( x\right)  = 2{x}^{3} + a{x}^{2} - {3x} - 2$ 在 $x = 1$ 处取得极小值，则函数 $f\left( x\right)$ 的极大值为\blkda{$-\dfrac{9}{8}$}.
    \begin{jieda}
         ${f}^{\prime }\left( x\right)  = 6{x}^{2} + {2ax} - 3$ ,由题意得 ${f}^{\prime }\left( 1\right)  = 6 + {2a} - 3 = 0$ ,解得 $a =  - \dfrac{3}{2}$，故 $f\left( x\right)  = 2{x}^{3} - \dfrac{3}{2}{x}^{2} - {3x} - 2,{f}^{\prime }\left( x\right)  = 6{x}^{2} - {3x} - 3 = 3\left( {{2x} + 1}\right) \left( {x - 1}\right)$ 
         \begin{dingautolist}{172}
             \item 当 $- \dfrac{1}{2} < x < 1$ 时, ${f}^{\prime }\left( x\right)  < 0,f\left( x\right)$ 单调递减
             \item 当 $x <  - \dfrac{1}{2}$ 或 $x > 1$ 时, ${f}^{\prime }\left( x\right)  > 0,f\left( x\right)$ 单调递增
         \end{dingautolist}
         故 $f\left( x\right)$ 在 $x =  - \dfrac{1}{2}$ 处取得极大值，故极大值为 $f\left( {-\dfrac{1}{2}}\right)  = 2 \times  \left( {-\dfrac{1}{8}}\right)  - \dfrac{3}{2} \times  \dfrac{1}{4} + \dfrac{3}{2} - 2 =  - \dfrac{9}{8}$ .
    \end{jieda}
    
    \item 已知函数 $f\left( x\right)  = {x}^{3} + \left( {a - 1}\right) {x}^{2} + \left( {a - 4}\right) x$ 是奇函数，则 $f\left( x\right)$ 的极小值是\blkda{$-2$}.
    \begin{jieda}
        $f(1) + f(-1) = 0$可得 $a = 1$，故 $f\left( x\right)  = {x}^{3} - {3x}$，经检验复合要求，${f}^{\prime }\left( x\right)  = 3{x}^{2} - 3$
        
        令 ${f}^{\prime }\left( x\right)  > 0$ ,即 $3{x}^{2} - 3 > 0$ ,解得 $x > 1$ 或 $x <  - 1$
        
        令 ${f}^{\prime }\left( x\right)  < 0$ ,即 $3{x}^{2} - 3 < 0$ ,解得 $- 1 < x < 1$
        
        则 $f\left( x\right)$ 在 $\left( {-\infty , - 1}\right) ,\left( {1, + \infty }\right)$ 上单调递增,在$(-1,1)$上单调递减,故 $f\left( x\right)$ 的极小值是 $f\left( 1\right)  =  - 2$
    \end{jieda}

    \item 若函数 $f\left( x\right)  = {x}^{3} - a{x}^{2} + {bx}\left( {a,b \in  \mathbf{R}}\right)$ 在 $x = 1$ 处取得极值 4,则 $a - b =$ \blkda{$-3$}.
    \begin{jieda}
        因为 $f\left( x\right)  = {x}^{3} - a{x}^{2} + {bx}$ 在 $x = 1$ 处取得极值 4，所以 ${f}^{\prime }\left( 1\right)  = 0$ 且 $f\left( 1\right)  = 4$.
        
        ${f}^{\prime }\left( x\right)  = 3{x}^{2} - {2ax} + b$ ,所以 ${f}^{\prime }\left( 1\right)  = 3 - {2a} + b = 0$，又 $f\left( 1\right)  = 1 - a + b = 4$，解得 $a = 6,b = 9$ ,经验证符合题意，所以 $a - b = 6 - 9 =  - 3$
    \end{jieda}

    \item 若函数 $f\left( x\right)  = {x}^{3} - a{x}^{2} + {2x} - 1$ 有极值，则实数 $a$ 的取值范围是\blkda[4]{$\left( {-\infty , - \sqrt{6}}\right)  \cup  \left( {\sqrt{6}, + \infty }\right)$}.
    \begin{jieda}
        由 $f\left( x\right)  = {x}^{3} - a{x}^{2} + {2x} - 1$，则 ${f}^{\prime }\left( x\right)  = 3{x}^{2} - {2ax} + 2$，由函数 $f\left( x\right)$ 有极值，即 ${f}^{\prime }\left( x\right)  = 3{x}^{2} - {2ax} + 2$ 有变号零点，所以 ${\left( -2a\right) }^{2} - 4 \cdot  3 \cdot  2 > 0$，故$a \in \left( {-\infty , - \sqrt{6}}\right)  \cup  \left( {\sqrt{6}, + \infty }\right)$
    \end{jieda}

    \item 若函数 $f\left( x\right)  = \dfrac{2}{3}{x}^{3} + a{x}^{2} + {8x} + 1$ 在区间 $\left( {1,3}\right)$ 上有极值，则 $a$ 的取值范围为\blkda{$(-5, -4)$}.
    \begin{jieda}
        ${f}^{\prime }\left( x\right)  = 2{x}^{2} + {2ax} + 8$，因函数 $f\left( x\right)  = \dfrac{2}{3}{x}^{3} + a{x}^{2} + {8x} + 1$ 在区间$(1,3)$上有极值，则方程 ${f}^{\prime }\left( x\right)  = 2{x}^{2} + {2ax} + 8 = 0$ 在区间$(1,3)$上有实根，故方程 $- a = x + \dfrac{4}{x}$ 在区间$(1,3)$上有实根，也即函数 $y =  - a$ 与 $g\left( x\right)  = x + \dfrac{4}{x}$ 在区间$(1,3)$上有交点.
        
        因 $g\left( x\right)  = x + \dfrac{4}{x}$ 在$(1,2)$上递减,在$(2,3)$上递增,且 $g\left( 1\right)  = 5 > g\left( 3\right)  = \dfrac{13}{3},g\left( 2\right)  = 4$，故 $4 \leq  g\left( x\right)  < 5$ ,即 $4 \leq   - a < 5$ ,解得 $- 5 < a \leq  -4$，当$a = -4$时，$f'(x) \geq 0$恒成立，故 $a$ 的取值范围为$(-5, -4)$.
    \end{jieda}

    \item 函数 $f\left( x\right)  = {x}^{3} + 3{x}^{2} + {3x} - a$ 的极值点的个数是\blkda{0}个.

    \item 已知函数 $f\left( x\right)  = \dfrac{1}{3}{a}^{2}{x}^{3} - \dfrac{5}{2}a{x}^{2} + {6x}$ 的极小值点为 1，则 $a =$ \blkda{3}.
    \begin{jieda}
        易知 ${f}^{\prime }\left( x\right)  = {a}^{2}{x}^{2} - {5ax} + 6 = \left( {{ax} - 2}\right) \left( {{ax} - 3}\right)$，由题意可知 $x = 1$ 是其一个变号零点，即 ${f}^{\prime }\left( 1\right)  = \left( {a - 2}\right) \left( {a - 3}\right)  = 0 \Longrightarrow  a = 2$ 或 $a = 3$
        \begin{dingautolist}{172}
            \item 当 $a = 2$ 时, ${f}^{\prime }\left( x\right)  = 2\left( {x - 1}\right) \left( {{2x} - 3}\right)$，此时 $x \in  \left( {1,\dfrac{3}{2}}\right)$ 时, ${f}^{\prime }\left( x\right)  < 0$ ,即 $f\left( x\right)$ 单调递减，$x \in  \left( {-\infty ,1}\right)$ , ${f}^{\prime }\left( x\right)  > 0$ ,即 $f\left( x\right)$ 单调递增,故在 $x = 1$ 处取得极大值,不符合题意
            \item 当 $a = 3$ 时, ${f}^{\prime }\left( x\right)  = 3\left( {{3x} - 2}\right) \left( {x - 1}\right)$，此时 $x \in  \left( {\dfrac{2}{3},1}\right)$ 时, ${f}^{\prime }\left( x\right)  < 0$ ,即 $f\left( x\right)$ 单调递减，$x \in  \left( {1, + \infty }\right)$ , ${f}^{\prime }\left( x\right)  > 0$ ,即 $f\left( x\right)$ 单调递增,故在 $x = 1$ 处取得极小值,符合题意
        \end{dingautolist}
        综上 $a = 3$ .
    \end{jieda}

    \item 设函数 $f\left( x\right)$ 在 $\mathbf{R}$ 上可导，其导函数为 ${f}^{\prime }\left( x\right)$ ，且函数 $y = \left( {1 - x}\right) {f}^{\prime }\left( x\right)$ 的图象如图所示，则下列结论中一定成立的是\dotfill{(\kh{A})}
    
    \begin{center}
    \includegraphics[width=0.2\textwidth]{bodyxb/src/bo_d6kkqnk601uc73e1cnh0_0_238_1754_332_309_0.jpg}
    \end{center}

    \xx{$f\left( x\right)$ 有极大值 $f\left( {-2}\right)$}{$f\left( x\right)$ 有极小值 $f\left( {-2}\right)$}{$f\left( x\right)$ 有极大值 $f\left( 1\right)$}{$f\left( x\right)$ 有极小值 $f\left( 1\right)$}
    
    \item 函数 $f\left( x\right)  = {x}^{3} + a{x}^{2} + {bx} + {a}^{2}$ 在 $x = 1$ 时有极值 10, 那么 $a$ 、 $b$ 的值为\blkda[4]{$a = -4, b = 11$}.
    \begin{jieda}
        因为函数 $f\left( x\right)  = {x}^{3} - a{x}^{2} - {bx} + {a}^{2}$ ,
        所以 ${f}^{\prime }\left( x\right)  = 3{x}^{2} - {2ax} - b$ ,
        因为函数 $f\left( x\right)  = {x}^{3} - a{x}^{2} - {bx} + {a}^{2}$ 在 $x = 1$ 时有极值 10,
        所以 $\left\{  \begin{array}{l} 1 - a - b + {a}^{2} = {10} \\  3 - {2a} - b = 0 \end{array}\right.$ ,
        即 $\left\{  \begin{array}{l}  - a - b + a^2 = 9 \\  {2a} + b = 3 \end{array}\right.$ ,
        解得 $\left\{  {\begin{array}{l} a =  - 4 \\  b = {11} \end{array}\text{ 或 }\left\{  {\begin{array}{l} a = 3 \\  b =  - 3 \end{array},}\right. }\right.$
        当 $\left\{  \begin{array}{l} a =  - 4 \\  b = {11} \end{array}\right.$ 时, ${f}^{\prime }\left( x\right)  = 3{x}^{2} + {8x} - {11} = \left( {{3x} + {11}}\right) \left( {x - 1}\right)$
        当 $x <  - \dfrac{11}{3}$ 或 $x > 1$ 时, ${f}^{\prime }\left( x\right)  > 0$ ,当 $- \dfrac{11}{3} < x < 1$ 时, ${f}^{\prime }\left( x\right)  < 0$ .
        所以 $f\left( x\right)$ 在 $x = 1$ 处取的极小值,符合题意.
        当 $\left\{  \begin{array}{l} a = 3 \\  b =  - 3 \end{array}\right.$ 时, ${f}^{\prime }\left( x\right)  = 3{x}^{2} - {6x} + 3 = 3{\left( x - 1\right) }^{2} \geq  0$
        所以 $f\left( x\right)$ 在 $x = 1$ 处无极值.
    \end{jieda}
    
    \item 函数 $y = f\left( x\right)$ 可导，“函数 $y = f\left( x\right)$ 在点 $\left( {{x}_{0},{y}_{0}}\right)$ 处的导数值为 0 ” 是 “函数 $y = f\left( x\right)$ 在点 $\left( {{x}_{0},{y}_{0}}\right)$ 处取极值” 的\dotfill{(\kh{B})}
    \xx{充分非必要条件}{必要非充分条件}{充要条件}{既非充分又非必要条件}
    
    \item 函数 $y = {x}^{3} - 3{x}^{2} - {9x}\left( {-{2 < x < 2}}\right)$ 有\dotfill{(\kh{C})}
    \xx{极大值 $5$ ，极小值 $-27$}{极大值 $5$ ，极小值 $-11$}{极大值 $5$ ，无极小值}{极小值 $-27$ ，无极大值}
    
    \item 函数 $f\left( x\right)$ 的定义域为开区间 $\left( {a,b}\right)$ ，导函数 ${f}^{\prime }\left( x\right)$ 在 $\left( {a,b}\right)$ 内的图象如图所示,则函数 $f\left( x\right)$ 在开区间 $\left( {a,b}\right)$ 内有极小值点\blkda{1}个.
        
    \begin{center}
    \includegraphics[width=0.3\textwidth]{bodyxb/src/bo_d6kkqnk601uc73e1cnh0_1_233_674_452_244_0.jpg}
    \end{center}
    
    \item 用导数求下列函数的极值:
    \begin{enumerate}
        \item $f\left( x\right)  = \dfrac{1}{2}{x}^{2} - {2x} + 4$
        \item $f\left( x\right)  = {x}^{3} - {x}^{2} - x + 1$
        \item $f\left( x\right)  = 3{x}^{4} - 8{x}^{3} + 6{x}^{2} + 1$ 
    \end{enumerate}
    \begin{jieda}
        \begin{enumerate}
            \item 极小值 $f\left( 2\right)  = 2$ ,无极大值
            \item 极大值 $f\left( {-\dfrac{1}{3}}\right)  = \dfrac{32}{27}$ ,极小值 $f\left( 1\right)  = 0$
            \item 极小值 $f\left( 0\right)  = 1$ ,无极大值
        \end{enumerate}
    \end{jieda}
    
    \item 求余弦函数 $y = \cos x$ 的单调区间和极值
    \begin{jieda}
        单调增区间为 $\left( {{2k\pi } - \pi ,{2k\pi }}\right) ,k \in  \mathbf{Z}$ ,单调减区间为 $\left( {{2k\pi },{2k\pi } + \pi }\right) ,k \in  \mathbf{Z}$ ,极大值$f\left( {2k\pi }\right)  = 1,k \in  \mathbf{Z}$ ,极小值 $f\left( {{2k\pi } + \pi }\right)  =  - 1,k \in  \mathbf{Z}$
    \end{jieda}
    
    \item 已知函数 $f\left( x\right)  = {{e}}^{x} - {3x} + {3a}\left( {a \in  \mathbf{R}}\right)$ ，求 $f\left( x\right)$ 的单调区间与极值.
    \begin{jieda}
        单调增区间为 $\left( {\ln 3, + \infty }\right)$ ,单调减区间为 $\left( {-\infty ,\ln 3}\right)$ ,极小值 $f\left( {\ln 3}\right)  = 3 - 3\ln 3 + {3a}$ ,无极大值.
    \end{jieda}
    
    \item 某函数图像如图所示，它在 $\left\lbrack  {a,b}\right\rbrack$ 上哪一点处取得最大值？它是极大值点吗? 在哪一点处取得最小值? 它是极小值吗?
    
    \begin{center}
    \includegraphics[width=0.6\textwidth]{bodyxb/src/bo_d6kkqnk601uc73e1cnh0_2_698_359_709_283_0.jpg}
    \end{center}
    \begin{jieda}
        $x = b$ 处取最大值,不是极大值点; $x = {x}_{4}$ 处取最大值,是极小值.
    \end{jieda}
    
    \item 为了支持企业复工复产，某地政府决定向当地企业发放补助款，其中对纳税额 $x$ (万元) 在 $\left\lbrack  {4,8}\right\rbrack$ 的小微企业做统一方案,方案要求同时具备下列两个条件:  \ding{172}  补助款 $f\left( x\right)$ (万元)随企业原纳税额 $x$ (万元)的增加而增加； \ding{173}  补助款不低于原纳税额的 50\%. 经测算政府决定采用函数模型 $f\left( x\right)  = \dfrac{x}{4} - \dfrac{m}{x} + 4$ 作为补助款发放方案.
    \begin{enumerate}
        \item 判断 $m = {12}$ 时是否满足条件，并说明理由
        \item 求同时满足条件 \ding{172}  \ding{173} 时 $m$ 的取值范围
    \end{enumerate}
    \begin{jieda}
        \begin{enumerate}
            \item 满足,对 $x \in  \left\lbrack  {4,8}\right\rbrack  ,{\left( f\left( x\right)  - \dfrac{1}{2}x\right) }^{\prime } =  - \dfrac{x}{4} - \dfrac{12}{x} + 4 \geq  0$
            \item $\left\lbrack  {-4,{12}}\right\rbrack$
        \end{enumerate}
    \end{jieda}
    
    \item 判断下列函数在 $\left( {-\infty , + \infty }\right)$ 上是否存在驻点,是否存在极值点,并说明理由: 
    \begin{enumerate*}
        \item $y = {x}^{n},n$ 为正奇数
        \item $y = {x}^{n}$ ， $n$ 为正偶数
    \end{enumerate*}
    \begin{jieda}
        \begin{enumerate}
            \item 当 $n = 1$ 时,无驻点,无极值点，当 $n \geq  3$ 时,驻点 $x = 0$ ,无极值点
            \item 驻点 $x = 0$ ,极小值点 $x = 0$
        \end{enumerate}
    \end{jieda}
    
    \item 已知函数 $y = {x}^{3} + {2m}{x}^{2} - {nx} + m$ 在 $x = 1$ 处有极值 0，求 $m + n$ 的值.
    \begin{jieda}
        $y' = 3x^2+4mx-n$，$y'|_{x = 1} = 3+4m-n = 0$，$1+2m-n+m = 0$，$m = -2, n = -5$，故$m+n = -7$
    \end{jieda}
    
    \item 已知 $f\left( x\right)  = a{x}^{2} + x - \ln x$ .
    \begin{enumerate}
        \item 若 $a = 1$ ，求函数 $y = f\left( x\right)$ 的极值
        \item 若 $y = f\left( x\right)$ 存在单调递增区间，求实数 $a$ 的取值范围
    \end{enumerate}
    \begin{jieda}
        \begin{enumerate}
            \item $a = 1$ 时, $f\left( x\right)  = {x}^{2} + x - \ln x$ 定义域为 $\left( {0, + \infty }\right)$ ,
            ${f}^{\prime }\left( x\right)  = {2x} + 1 - \dfrac{1}{x} = \dfrac{2{x}^{2} + x - 1}{x} = \dfrac{\left( {{2x} - 1}\right) \left( {x + 1}\right) }{x},x > 0$
            令 ${f}^{\prime }\left( x\right)  = 0$ ,解得 $x = \dfrac{1}{2}$ .
            当 $x \in  \left( {0,\dfrac{1}{2}}\right)$ 时, ${f}^{\prime }\left( x\right)  < 0,f\left( x\right)$ 单调递减;
            当 $x \in  \left( {\dfrac{1}{2}, + \infty }\right)$ 时, ${f}^{\prime }\left( x\right)  > 0,f\left( x\right)$ 单调递增.
            所以在 $x = \dfrac{1}{2}$ 时,函数 $f\left( x\right)$ 取得极小值 $f\left( \dfrac{1}{2}\right)  = \dfrac{3}{4} + \ln 2;f\left( x\right)$ 无极大值.
            \item ${f}^{\prime }\left( x\right)  = {2ax} + 1 - \dfrac{1}{x} = \dfrac{{2a}{x}^{2} + x - 1}{x},\left( {x > 0}\right)$
            因为 $y = f\left( x\right)$ 存在单调递增区间,所以 ${2a}{x}^{2} + x - 1 > 0$ 在 $\left( {0, + \infty }\right)$ 上有解.
            由 ${2a}{x}^{2} + x - 1 > 0,x \in  \left( {0, + \infty }\right)$ 得 ${2a} > \dfrac{1}{{x}^{2}} - \dfrac{1}{x} = {\left( \dfrac{1}{x} - \dfrac{1}{2}\right) }^{2} - \dfrac{1}{4}$ ,
            $\because {\left( \dfrac{1}{x} - \dfrac{1}{2}\right) }^{2} - \dfrac{1}{4} \geq   - \dfrac{1}{4}$ ,当且仅当 $x = 2$ 时取等号.
            $\therefore {2a} >  - \dfrac{1}{4}$ ,即 $a >  - \dfrac{1}{8}$ ,
            所以实数 $a$ 的取值范围是 $\left( {-\dfrac{1}{8}, + \infty }\right)$ .
        \end{enumerate}
    \end{jieda}
    \item 已知函数 $f\left( x\right)  = 3\ln x - x + \dfrac{a}{x}$ .
    \begin{enumerate}
        \item 若 $a = 2$ ，求 $f\left( x\right)$ 的单调区间
        \item 若 $f\left( x\right)$ 没有极值点，求 $a$ 的取值范围
    \end{enumerate}
    \begin{jieda}
        \begin{enumerate}
            \item 若 $a = 2$ ,则 $f\left( x\right)  = 3\ln x - x + \dfrac{2}{x}$ ,函数定义域为 $\left( {0, + \infty }\right)$ ,
            ${f}^{\prime }\left( x\right)  = \dfrac{3}{x} - 1 - \dfrac{2}{{x}^{2}} =  - \dfrac{\left( {x - 1}\right) \left( {x - 2}\right) }{{x}^{2}}.$
            当 $0 < x < 1$ 时, ${f}^{\prime }\left( x\right)  < 0$ ；
            当 $1 < x < 2$ 时, ${f}^{\prime }\left( x\right)  > 0$ ；
            当 $x > 2$ 时, ${f}^{\prime }\left( x\right)  < 0$ ,
            故 $f\left( x\right)$ 的单调递减区间是$(0,1)$和 $\left( {2, + \infty }\right)$ ,单调递增区间是$(1,2)$.
            \item ${f}^{\prime }\left( x\right)  = \dfrac{3}{x} - 1 - \dfrac{a}{{x}^{2}} =  - \dfrac{{x}^{2} - {3x} + a}{{x}^{2}}$ ,
            函数 $h\left( x\right)  = {x}^{2} - {3x} + a$ ,当 $\Delta  = 9 - {4a} \leq  0$ ,即 $a \geq  \dfrac{9}{4}$ 时, $h\left( x\right)  \geq  0$ 恒成立，则有 ${f}^{\prime }\left( x\right)  \leq  0,f\left( x\right)$ 单调递减，此时 $f\left( x\right)$ 没有极值点,符合题意.
            当 $a < \dfrac{9}{4}$ 时,方程 ${x}^{2} - {3x} + a = 0$ 有两个实数根 ${x}_{1},{x}_{2}$ ,不妨设 ${x}_{1} < {x}_{2}$ ,
            则 ${x}_{1} + {x}_{2} = 3 > 0,{x}_{1}{x}_{2} = a$ .
            当 $0 < a < \dfrac{9}{4}$ 时, $0 < {x}_{1} < {x}_{2}$ ,此时 $f\left( x\right)$ 在区间 $\left( {0,{x}_{1}}\right) ,\left( {{x}_{2}, + \infty }\right)$ 上单调递减,
            在区间 $\left( {{x}_{1},{x}_{2}}\right)$ 上单调递增,所以 ${x}_{1}$ 是 $f\left( x\right)$ 的极小值点, ${x}_{2}$ 是 $f\left( x\right)$ 的极大值点,不符合题意;
            同理可知,当 $a \leq  0$ 时, $f\left( x\right)$ 在区间 $\left( {0,{x}_{2}}\right)$ 上单调递增, $\left( {{x}_{2}, + \infty }\right)$ 上单调递减, ${x}_{2}$ 是 $f\left( x\right)$ 的
            极大值点, 不符合题意.
            综上, $a$ 的取值范围是 $\left\lbrack  {\dfrac{9}{4}, + \infty }\right)$
        \end{enumerate}
    \end{jieda}
\end{enumerate}

C组：
\begin{enumerate}
    \item 设函数 $f\left( x\right)  = \sin x - x\cos x$ 在区间 $\left( {0,{10\pi }}\right)$ 内的极值点分别为 ${x}_{1},{x}_{2},\cdots ,{x}_{n}$ ，则 $f\left( {x}_{1}\right)  + f\left( {x}_{2}\right)  + \cdots  + f\left( {x}_{n}\right)  =$ \blkda{$5\pi$}.
    \begin{jieda}
        由 $f\left( x\right)  = \sin x - x\cos x$ 得 ${f}^{\prime }\left( x\right)  = x\sin x,x \in  \left( {0,{10\pi }}\right)$ ,
        令 ${f}^{\prime }\left( x\right)  = 0$ ,得 $x = \pi ,{2\pi },{3\pi },,{9\pi }$ ,
        由 $f\left( x\right)  = \sin x - x\cos x$ 为连续不断的曲线及极值点的概念可知:
        ${x}_{1} = \pi ,{x}_{2} = {2\pi },{x}_{3} = {3\pi },\cdots ,{x}_{9} = {9\pi }$ ,
        则 $f\left( {x}_{1}\right)  + f\left( {x}_{2}\right)  + f\left( {x}_{3}\right)  + \cdots  + f\left( {x}_{9}\right)  =  - \pi \cos \pi  - {2\pi }\cos {2\pi } - {3\pi }\cos {3\pi } - \cdots  - {9\pi }\cos {9\pi } = \pi  - {2\pi } + {3\pi } - \cdots  + {9\pi } = {5\pi }$ .
    \end{jieda}
    
    \item 函数 $f\left( x\right)  = 2 - {{e}}^{x} + m{x}^{2}$ 有两个极值点，则实数 $m$ 的取值范围是\blkda{$\left(\dfrac{e}{2}, +\infty\right)$}.
    \begin{jieda}
        由 $f\left( x\right)  = 2 - {{e}}^{x} + m{x}^{2}$ 得 ${f}^{\prime }\left( x\right)  =  - {{e}}^{x} + {2mx}$，因为函数 $f\left( x\right)  = 2 - {{e}}^{x} + m{x}^{2}$ 有两个极值点,所以 ${f}^{\prime }\left( x\right)  =  - {{e}}^{x} + {2mx}$ 有两个异号零点，即 ${{e}}^{x} = {2mx}$ 有两个不同的根,显然 $x \neq  0$ ,则 $\dfrac{{{e}}^{x}}{x} = {2m}$ 有两个不同的根,令 $g\left( x\right)  = \dfrac{{{e}}^{x}}{x}$ ，则 $y = {2m}$ 与 $g\left( x\right)  = \dfrac{{{e}}^{x}}{x}$ 的图象有两个不同的交点
        
        ${g}^{\prime }\left( x\right)  = \dfrac{{{e}}^{x}\left( {x - 1}\right) }{{x}^{2}}$ ,当 $x < 0$ 和 $0 < x < 1$ 时, ${g}^{\prime }\left( x\right)  < 0,g\left( x\right)$ 单调递减，当 $x > 1$ 时, ${g}^{\prime }\left( x\right)  > 0$ , $g\left( x\right)$ 单调递增,在 $x = 1$ 时,取得极小值 $g\left( 1\right)  = {e}$ ,作出图象如图：
            \begin{center}
                \includegraphics[width=0.1\linewidth]{bodyxb/src/jizhiC2.png}
            \end{center}
        由图可知 ${2m} > {e}$ ,所以 $m > \dfrac{{e}}{2}$ ,所以 $m$ 的取值范围是 $\left( {\dfrac{{e}}{2}, + \infty }\right)$ .
    \end{jieda}

    \item 已知 $\omega  > 0$ 且函数 $f\left( x\right)  = \sin \left( {{\omega x} + \dfrac{\pi }{4}}\right)$ 的图象恰有三条对称轴在 $\left( {0,\pi }\right)$ 上，则其在 $\left( {0,\pi }\right)$ 上共有\blkda{2}个极大值点.
    \begin{jieda}
        原命题即$g(t) = sint$在$\left(\dfrac{\pi}{4}, \omega \pi + \dfrac{\pi}{4}\right)$上有三条对称轴，根据正弦图像可知有两个极大值点.
    \end{jieda}
\end{enumerate}
\clearpage
\subsection*{导数的应用(3)——函数的最值}

A组：
\begin{enumerate}
    \item 函数 $f\left( x\right)  = 2{x}^{3} - 2{x}^{2}$ 在区间 $\left\lbrack  {-1,2}\right\rbrack$ 上的最大值是\blkda{8}.

    \item 已知函数 $f\left( x\right)  = \dfrac{1}{2}x - \sin x$ ，则 $f\left( x\right)$ 在 $\left\lbrack  {0,\pi }\right\rbrack$ 上的值域为\blkda[3]{$\left[\dfrac{\pi}{6}-\dfrac{\sqrt{3}}{2}, \dfrac{\pi}{2}\right]$}.

    \item 函数图象连续的函数 $y = f\left( x\right)$ 在区间 $\left\lbrack  {a,b}\right\rbrack$ 上\dotfill{(\kh{C})}
    \xx{一定存在极小值}{一定存在极大值}{一定存在最大值}{极小值一定比极大值小}
    
    \item 下列有关函数的极值与最值的命题中，为真命题的是\dotfill{(\kh{B})}
    \xx{函数的最大值一定不是这个函数的极大值}{函数的极大值可以小于这个函数的极小值}{函数在某一闭区间上的极小值就是函数的最小值}{函数在开区间上不存在极大值和最大值}

    \item 求下列函数在给定区间上的最大值与最小值:
    \begin{enumerate}
        \item $f\left( x\right)  = x - {x}^{3},x \in  \left\lbrack  {0,2}\right\rbrack$ ; 
        
        \item $f\left( x\right)  =  - 3{x}^{4} + 6{x}^{2} - 1,x \in  \left\lbrack  {-2,2}\right\rbrack$

        \item $y = {x}^{2}{e}^{x},x \in  \left\lbrack  {-3,1}\right\rbrack$
    \end{enumerate}
    \begin{jieda}
        \begin{enumerate}
            \item 最大值为 $\dfrac{2\sqrt{3}}{9}$，最小值为 $-6$.
            \item 最大值为 $2$，最小值为 $-25$
            \item 最小值为 $0$，最大值$e$
        \end{enumerate}
    \end{jieda}
\end{enumerate}

B组：
\begin{enumerate}
    \item 函数 $f\left( x\right)  = \dfrac{1}{2}{x}^{2} - \left( {{e} - 1}\right) x - {{eln}}x$ 的最小值为\blkda{$- \dfrac{1}{2}{{e}}^{2}$}.
    \begin{jieda}
        易知 ${f}^{\prime }\left( x\right)  = x - \left( {{e} - 1}\right)  - \dfrac{{e}}{x} = \dfrac{\left( {x + 1}\right) \left( {x - {e}}\right) }{x}\left( {x > 0}\right)$，所以 $x > {e}$ 时, ${f}^{\prime }\left( x\right)  > 0$ ,即此时函数单调递增，$x \in  \left( {0,{e}}\right)$ 时, ${f}^{\prime }\left( x\right)  < 0$ ,即此时函数单调递减，所以 $f\left( x\right)  \geq  f\left( {e}\right)  = \dfrac{1}{2}{{e}}^{2} - \left( {{{e}}^{2} - {e}}\right)  - {e} =  - \dfrac{1}{2}{{e}}^{2}$，即该函数的最小值为 $- \dfrac{1}{2}{{e}}^{2}$
    \end{jieda}

    \item 已知函数 $f\left( x\right)  = \dfrac{{x}^{3}}{{{e}}^{x}}$ ，则 $f\left( x\right)$ 的最大值为\blkda{$\dfrac{27}{{{e}}^{3}}$}.
    \begin{jieda}
        函数 $f\left( x\right)  = \dfrac{{x}^{3}}{{{e}}^{x}}$ 的定义域为 $\mathbf{R}$ ,求导得 ${f}^{\prime }\left( x\right)  = \dfrac{{x}^{2}\left( {3 - x}\right) }{{{e}}^{x}}$
        \begin{dingautolist}{172}
            \item 当 $x < 3$ 时, ${f}^{\prime }\left( x\right)  > 0$ 
            \item 当 $x > 3$ 时, ${f}^{\prime }\left( x\right)  < 0$ 
        \end{dingautolist}
        
        函数 $f\left( x\right)$ 在 $\left( {-\infty ,3}\right)$ 上单调递增,在 $\left( {3, + \infty }\right)$ 上单调递减，所以当 $x = 3$ 时, $f\left( x\right)$ 取得最大值 $f\left( 3\right)  = \dfrac{27}{{{e}}^{3}}$ .
    \end{jieda}

    \item 函数 $y =  - {x}^{4} + {x}^{3} + 1$ 在区间 $\left\lbrack  {-1,1}\right\rbrack$ 上的最大值与最小值之和为\blkda{$\dfrac{27}{256}$}.

    \item 将一个边长为 $3{m}$ 的正方形铁片的四角截去四个边长相等的小正方形，做成一个无盖方盒，则方盒的容积的最大值为\blkda{2} ${{m}}^{3}$.
    \begin{jieda}
        依题意方盒的底面边长为 $\left( {3 - {2x}}\right) {m}$ 的正方形,高为 $x{m}$，则 $\left\{  \begin{array}{l} 3 - {2x} > 0 \\  x > 0 \end{array}\right.$ ,即 $0 < x < \dfrac{3}{2}$，所以无盖方盒的容积为 $V\left( x\right)  = x{\left(3 - 2x\right) }^{2},x \in  \left( {0,\dfrac{3}{2}}\right)$，则 ${V}^{\prime }\left( x\right)  = {12}{x}^{2} - {24x} + 9$ 
        
        令 ${V}^{\prime }\left( x\right)  = {12}{x}^{2} - {24x} + 9 > 0$ ,解得 $x < \dfrac{1}{2}$ 或 $x > \dfrac{3}{2}$
        
        令 ${V}^{\prime }\left( x\right)  = {12}{x}^{2} - {24x} + 9 < 0$ ,解得 $\dfrac{1}{2} < x < \dfrac{3}{2}$ .
        
        因为函数 $V\left( {x}\right)$ 的定义域为 $x \in  \left( {0,\dfrac{3}{2}}\right)$，所以函数 $V\left( {x}\right)$ 在 $\left( {0,\dfrac{1}{2}}\right)$ 上单调递增,在 $\left( {\dfrac{1}{2},\dfrac{3}{2}}\right)$ 上单调递减，所以 $V\left( x\right)$ 在 $x = \dfrac{1}{2}$ 处取得极大值即最大值,所以 $V{\left( x\right) }_{\max } = V\left( \dfrac{1}{2}\right)  = 2$，即该方盒容积最大为 $2{{m}}^{3}$ .
    \end{jieda}

    \item 若曲线 $f\left( x\right)  = a\ln x - x - 2$ 在 $x = 1$ 处有最值，则实数 $a$ 的值为\blkda{1}.
    \begin{jieda}
        曲线 $f\left( x\right)  = a\ln x - x - 2$ ,定义域为 $\left( {0, + \infty }\right)$ ,所以 ${f}^{\prime }\left( x\right)  = \dfrac{a}{x} - 1 = \dfrac{a - x}{x}$，当 $a \leq  0$ 时, ${f}^{\prime }\left( x\right)  < 0$ 所以 $f\left( x\right)$ 单调递减,无最大值不合题意;
        
        当 $a > 0$ 时, $x > a,{f}^{\prime }\left( x\right)  < 0,f\left( x\right)$ 单调递减, $0 < x <a, {f}^{\prime }\left( x\right) >  0,f\left( x\right)$ 单调递增,所以 $x = a$ 时函数取最大值，因为函数在 $x = 1$ 处有最值,所以 $a = 1$
    \end{jieda}

    \item 已知函数 $f\left( x\right)  = \left\{  \begin{array}{l} {{e}}^{x},x > 0 \\  {x}^{3} - {3x} + a,x \leq  0 \end{array}\right.$ 的值域为 $\mathbf{R}$ ，则实数 $a$ 的取值范围为\blkda{$[-1, +\infty)$}.
    \begin{jieda}
        当 $x > 0$ 时, $f\left( x\right)  = {e}^{x}$ 的值域为 $\left( {1, + \infty }\right)$

        故应有当 $x \leq  0$ 时, $( - \infty ,1\rbrack$ 是 $f\left( x\right)  = {x}^{3} - {3x} + a$ 的值域的子集，又 ${f}^{\prime }\left( x\right)  = 3{x}^{2} - 3 = 3\left( {x + 1}\right) \left( {x - 1}\right)$ ,令 ${f}^{\prime }\left( x\right)  = 0,\therefore x =  - 1$ 或 $x = 1$ (舍去)
        
        当 $x <  - 1$ 时, ${f}^{\prime }\left( x\right)  > 0$ , $f\left( x\right)$ 单调递增；当 $- 1 < x \leq  0$ 时, ${f}^{\prime }\left( x\right)  < 0,f\left( x\right)$ 单调递减
        
        当 $x =  - 1$ 时, $f\left( x\right)$ 取得极大值 $f\left( {-1}\right)  = 2 + a$，故 $f\left( x\right)  = {x}^{3} - {3x} + a\left( {x \leq  0}\right)$ 的值域为 $( - \infty ,2 + a\rbrack$
        
        故$2 + a \geq  1$，故$a \geq   - 1$ .
    \end{jieda}

    \item 已知 ${f}^{\prime }\left( x\right)$ 是函数 $f\left( x\right)  = {{e}}^{x} + {ax} + c\left( {a,c \in  \mathbf{R}}\right)$ 的导函数，若 ${f}^{\prime }\left( 0\right)  = 3$ ，且 $f\left( x\right)$ 在 $\left\lbrack  {0,1}\right\rbrack$ 上的最大值为 5，则 $c$ 的值为\blkda{$3-e$}.
    \begin{jieda}
        因为 $f\left( x\right)  = {{e}}^{x} + {ax} + c\left( {a,c \in  \mathbf{R}}\right)$ ,所以 ${f}^{\prime }\left( x\right)  = {{e}}^{x} + a$，所以 ${f}^{\prime }\left( 0\right)  = {{e}}^{0} + a = 3$ ,解得 $a = 2$，所以 $f\left( x\right)  = {{e}}^{x} + {2x} + c$ ,则 ${f}^{\prime }\left( x\right)  = {{e}}^{x} + 2$，所以当 $x \in  \left\lbrack  {0,1}\right\rbrack$ 时 ${f}^{\prime }\left( x\right)  > 0$ ,所以 $f\left( x\right)$ 在 $\left\lbrack  {0,1}\right\rbrack$ 上单调递增，所以 $f\left( 1\right)  = {e} + 2 + c = 5$ ,解得 $c = 3 - {e}$
    \end{jieda}

    \item 若函数 $f\left( x\right)  =  -{x}^{3} + \dfrac{3}{2}{x}^{2}$ 在区间 $\left( {{2a} - 5,a + 1}\right)$ 内有最小值，则实数 $a$ 的取值范围是\blkda{$\left(-1, \dfrac{1}{2}\right]$}.
    \begin{jieda}
        由题可知: ${f}^{\prime }\left( x\right)  =  - 3{x}^{2} + {3x} =  - {3x}\left( {x - 1}\right)$，令 ${f}^{\prime }\left( x\right)  > 0$，则 $0 < x < 1$ ; 令 ${f}^{\prime }\left( x\right)  < 0$ ,则 $x < 0$ 或 $x > 1$，所以函数 $f\left( x\right)$ 在$(0,1)$单调递增,在 $\left( {-\infty ,0}\right) ,\left( {1, + \infty }\right)$ 单调递减.
        
        极小值为 $f\left( 0\right)  = 0$ ,令 $f\left( x\right)  = 0$ ,所以 $x = 0$ 或 $\dfrac{3}{2}$，又函数 $f\left( x\right)  =  - {x}^{3} + \dfrac{3}{2}{x}^{2}$ 在区间 $\left( {{2a} - 5,a + 1}\right)$ 内有最小值，所以 $\left\{  \begin{array}{l} {2a} - 5 < 0 \\  0 < a + 1 \leq  \dfrac{3}{2} \end{array}\right.  \Longrightarrow   - 1 < a \leq  \dfrac{1}{2}$.
    \end{jieda}

    \item 已知函数 $f\left( x\right)  = \ln x - {ax} - 2$ 在区间 $\left( {1,2}\right)$ 上存在最大值,则实数 $a$ 的取值范围为\blkda{$\left( {\dfrac{1}{2},1}\right)$}.
    \begin{jieda}
        $f\left( x\right)  = \ln x - {ax} - 2 \Longrightarrow  {f}^{\prime }\left( x\right)  = \dfrac{1}{x} - a = \dfrac{1 - {ax}}{x}$
        \begin{dingautolist}{172}
            \item 当 $a \leq  0$ 时, $f\left( x\right)$ 在 $\left( {0, + \infty }\right)$ 上严格单调递增,不符合题意
            \item 当 $a > 0$ 时,令 ${f}^{\prime }\left( x\right)  > 0 \Longrightarrow  0 < x < \dfrac{1}{a};{f}^{\prime }\left( x\right)  < 0 \Longrightarrow  x > \dfrac{1}{a}$
        \end{dingautolist}
        所以 $f\left( x\right)$ 在 $\left( {0,\dfrac{1}{a}}\right)$ 上严格单调递增,在 $\left( {\dfrac{1}{a}, + \infty }\right)$ 上严格单调递减，所以 $f\left( x\right)$ 在 $x = \dfrac{1}{a}$ 处取得极大值.
        
        因为函数 $f\left( x\right)$ 在区间$(1,2)$上存在最大值. 所以 $1 < \dfrac{1}{a} < 2 \Longrightarrow  a \in  \left( {\dfrac{1}{2},1}\right)$ 
    \end{jieda}

    \item 把三根结实且等长的树干一端用藤条捆扎起来，另一端立在地面不同的位置得到一个正三棱锥架构, 再用树枝、杂草等物覆盖形成侧面, 可在野外搭建起一个三棱锥形状的简易帐篷, 能起到遮风挡雨的作用, 设三根树干的长度都为 6 , 当帐篷的容积最大时(不计损耗)，其高度为\dotfill{(\kh{C})}
    \xx{2}{3}{$2\sqrt{3}$}{$3\sqrt{2}$}
    \begin{jieda}
        设帐篷高度为 $h\left( {0 < h < 6}\right)$，则底面正三角形的外接圆半径 $r = \sqrt{{36} - {h}^{2}}$，易知底面边长 $a = \sqrt{3}r = \sqrt{3}\sqrt{{36} - {h}^{2}}$，底面面积为 $S = \dfrac{\sqrt{3}}{4}{a}^{2} = \dfrac{3\sqrt{3}}{4}\left( {{36} - {h}^{2}}\right)$，帐篷容积 $V\left( h\right)  = \dfrac{1}{3} \times  \dfrac{3\sqrt{3}}{4}\left( {{36} - {h}^{2}}\right)  \cdot  h = \dfrac{\sqrt{3}}{4}\left( {{36h} - {h}^{3}}\right)$，则 ${V}^{\prime }\left( h\right)  = \dfrac{\sqrt{3}}{4}\left( {{36} - 3{h}^{2}}\right)$
        
        令 ${V}^{\prime }\left( h\right)  = 0$ 得 $h = 2\sqrt{3}$
        
        当 $h \in  \left( {0,2\sqrt{3}}\right)$ 时, ${V}^{\prime }\left( h\right)  > 0,V\left( h\right)$ 单调递增
        
        当 $h \in  \left( {2\sqrt{3},6}\right)$ 时, ${V}^{\prime }\left( h\right)  < 0,V\left( h\right)$ 单调递减
        
        所以 $V\left( h\right)$ 在 $h = 2\sqrt{3}$ 时取得最大值.
    \end{jieda}
    
    \item 已知函数 $y = f\left( x\right)$ 定义域为 $\left( {-1,2}\right)$ ，且在该区间上连续，在 $\left( {-1,2}\right)$ 上函数 $y = f\left( x\right)$ 有唯一的极大值 $f\left( 1\right)$ ，则下列说法正确的是\dotfill{(\kh{D})}
    \xx{函数 $f\left( x\right)$ 有最大值 $f\left( 1\right)$}{函数 $f\left( x\right)$ 有最大值,但不一定是 $f\left( 1\right)$}{函数 $f\left( x\right)$ 所显小值也可能是 $f\left( 1\right)$}{函数 $f\left( x\right)$ 不一定有最大值}
    
    \item 已知函数 $f\left( x\right)  = \dfrac{1}{3}{x}^{3} - {4x} + 4$
    \begin{enumerate}
        \item 求 $f\left( x\right)$ 的单调区间和极值
        \item 当$x \in [0, 3]$时，求$f(x)$的最值
    \end{enumerate}
    \begin{jieda}
        \begin{enumerate}
            \item 由 $f\left( x\right)  = \dfrac{1}{3}{x}^{3} - {4x} + 4$ ,得 ${f}^{\prime }\left( x\right)  = {x}^{2} - 4$，令 ${f}^{\prime }\left( x\right)  = 0$ ,解得 $x =  \pm  2$
            
            当 $x \in  \left( {-\infty , - 2}\right)$ 时 ${f}^{\prime }\left( x\right)  > 0$ ,所以函数 $f\left( x\right)$ 在 $\left( {-\infty , - 2}\right)$ 上单调递增
            
            当 $x \in  \left( {2, + \infty }\right)$ 时 ${f}^{\prime }\left( x\right)  > 0$ ,所以函数 $f\left( x\right)$ 在 $\left( {2, + \infty }\right)$ 上单调递增
            
            当 $x \in  \left( {-2,2}\right)$ 时 ${f}^{\prime }\left( x\right)  < 0$ ,所以函数 $f\left( x\right)$ 在(-2,2)上单调递减
            
            所以当 $x =  - 2$ 时,函数 $f\left( x\right)$ 有极大值为 $f\left( {-2}\right)  = \dfrac{1}{3} \times  {\left( -2\right) }^{3} - 4 \times  \left( {-2}\right)  + 4 = \dfrac{28}{3}$ 
            
            当 $x = 2$ 时,函数 $f\left( x\right)$ 有极小值为 $f\left( 2\right)  = \dfrac{1}{3} \times  {2}^{3} - 4 \times  2 + 4 =  - \dfrac{4}{3}$ .
            
            综上, $f\left( x\right)$ 的单调递增区间为 $\left( {-\infty , - 2}\right) ,\left( {2, + \infty }\right)$ ,单调递减区间为$(-2,2)$，故函数 $f\left( x\right)$ 有极大值为 $\dfrac{28}{3}$ ,有极小值为 $- \dfrac{4}{3}$ 
            
            \item 由(1)得 $f\left( x\right)$ 在 $\left\lbrack  {0,2}\right\rbrack$ 上单调递减,在 $(2,3\rbrack$ 上单调递增,
            又 $f\left( 0\right)  = 4,f\left( 3\right)  = \dfrac{1}{3} \times  {3}^{3} - 4 \times  3 + 4 = 1$ ,
            又函数 $f\left( x\right)$ 的极小值为 $f\left( 2\right)  =  - \dfrac{4}{3}$，所以当 $x \in  \left\lbrack  {0,3}\right\rbrack$ 时,函数 $f\left( x\right)$ 的最大值为 4 ,最小值为 $- \dfrac{4}{3}$.
        \end{enumerate}
    \end{jieda}
    
    \item 已知 ${f}\left( {x}\right)  = \sqrt{{x}},{g}\left( {x}\right)  = {k}{{x}}^{2}$ 
    \begin{enumerate}
        \item 求曲线 ${y} = {f}\left( {x}\right)$ 在点 $\left( {{4},{2}}\right)$ 处的切线方程
        \item 若曲线 $y = g\left( x\right)$ 经过点 $\left( {4,2}\right)$ ,求它与(1)中切线的另一个交点
    \end{enumerate}
    \begin{jieda}
        \begin{enumerate}
            \item $y = \dfrac{x}{4}+1$
            \item $\left(-2, \dfrac{1}{2}\right)$
        \end{enumerate}
    \end{jieda}
    
    \item 设 $x =  - 1$ 是函数 $f\left( x\right)  =  - {x}^{3} + a{x}^{2} + {bx} + c$ 的一个驻点,曲线 $y = f\left( x\right)$ 在 $x = 2$ 处的切线斜率为 9.
    \begin{enumerate}
        \item 求 $f\left( x\right)$ 的单调区间
        \item 若 $f\left( x\right)$ 在闭区间 $\left\lbrack  {-2,2}\right\rbrack$ 上的最大值为 20,求 $c$ 的值.
    \end{enumerate}
    \begin{jieda}
        \begin{enumerate}
            \item ${f}^{\prime }\left( x\right)  =  - 3{x}^{2} + {2ax} + b$ ,由已知得 $\left\{  \begin{array}{l} {f}^{\prime }\left( {-1}\right)  = 0 \\  {f}^{\prime }\left( 2\right)  = 9 \end{array}\right.$ ,得 $\left\{  \begin{array}{l}  - 3 - {2a} + b = 0 \\   - {12} + {4a} + b = 9 \end{array}\right.$
            解得 $a = 3,b = 9$ . 经验证可知符合题意,
            于是 ${f}^{\prime }\left( x\right)  =  - 3{x}^{2} + {6x} + 9 =  - 3\left( {x + 1}\right) \left( {x - 3}\right)$ ,
            由 ${f}^{\prime }\left( x\right)  > 0$ ,得 $- 1 < x < 3$ ,由 ${f}^{\prime }\left( x\right)  < 0$ ,得 $x <  - 1$ 或 $x > 3$ ,
            所以 $f\left( x\right)$ 的单调递增区间是$(-1,3)$,单调递减区间是 $\left( {-\infty , - 1}\right)$ 和 $\left( {3, + \infty }\right)$ .
            \item 由(1)知 $f\left( x\right)  =  - {x}^{3} + {3{x}^{2}} + {9x} + c$ ,因为 $f\left( x\right)$ 在区间 $\lbrack  - 2, - 1)$ 上是单调递减函数,在 $( - 1,2\rbrack$ 上是单调递增函数,又 $f\left( {-2}\right)  =$
            $2 + c < f\left( 2\right)  = {22} + c$，所以其最大值为 $f\left( 2\right)  = {22} + c = {20}$ ,解得 $c =  - 2$ .
        \end{enumerate}
    \end{jieda}
    
    \item 已知函数 $f\left( x\right)  = x{\log }_{2}x + \left( {1 - x}\right) {\log }_{2}\left( {1 - x}\right) ,x \in  \left( {0,1}\right)$ ,求函数 $y = f\left( x\right)$ 的最大值和最小值.
    \begin{jieda}
        最小值$-1$，无最大值.
    \end{jieda}
    
    \item 已知函数 $f\left( x\right)  = 2\cos {2x} - 4\sin x$ ,求函数 $y = f\left( x\right)$ 的最大值和最小值.
    \begin{jieda}
        最大值为$3$，最小值为$-6$.
    \end{jieda}
    
    \item 利用导数判断函数 $y = \dfrac{1}{\cos x},x \in  \left( {-\dfrac{\pi }{2},\dfrac{\pi }{2}}\right)$ 的单调性,并求出极值.
    \begin{jieda}
        单调减区间 $\left( {-\dfrac{\pi }{2},0}\right)$ ,单调增区间 $\left( {0,\dfrac{\pi }{2}}\right)$ ,极小值 $f\left( 0\right)  = 1$ ,无极大值
    \end{jieda}

    \item 已知 $a\text{ 、 }\beta$ 是方程 $4{x}^{2} - {4tx} - 1 = 0\left( {t \in  \mathbf{R}}\right)$ 的两个不等实根,函数 $f\left( x\right)  = \dfrac{{2x} - t}{{x}^{2} + 1}$ 的定义域为 $\left\lbrack  {a,\beta }\right\rbrack$ ,记函数 $f\left( x\right)$ 的最大值为 $\operatorname{MAX}f\left( x\right)$ ,最小值为 $\operatorname{MIN}f\left( x\right)$ , $g\left( t\right)  = \operatorname{MAX}f\left( x\right)  - \operatorname{MIN}f\left( x\right)$ ,求 $g\left( t\right)$ 的表达式.
    \begin{jieda}
        由于 $\alpha \text{、}\beta$ 是方程 ${4x2} - {4tx} - 1 = 0\left( {t \in  \mathbf{R}}\right)$ 的两个不等实根,
        所以 $\Delta  = {16}{t}^{2} + {16} > 0$ ,且 $\alpha  + \beta  = t,{\alpha \beta } =  - \dfrac{1}{4}$ .
        又函数 $f\left( x\right)  = \dfrac{{2x} - t}{{x}^{2} + 1}$ 的定义域为 $\left\lbrack  {\alpha ,\beta }\right\rbrack$ ,则 $\alpha  \leq  x \leq  \beta$ ,即 $\left( {x - \alpha }\right) \left( {x - \beta }\right)  \leq  0$ .
        对 $f\left( x\right)$ 求导数,得 ${f}^{\prime }\left( x\right)  = \dfrac{2\left( {{x}^{2} + 1}\right)  - \left( {{2x} - t}\right)  \cdot  {2x}}{{\left( {x}^{2} + 1\right) }^{2}} = \dfrac{-2{x}^{2} + {2tx} + 2}{{\left( {x}^{2} + 1\right) }^{2}} = \dfrac{-4{x}^{2} + {4tx} + 4}{2{\left( {x}^{2} + 1\right) }^{2}} = \dfrac{-4{x}^{2} + 4\left( {\alpha  + \beta }\right) x - {4\alpha \beta } + 3}{2{\left( {x}^{2} + 1\right) }^{2}} = \dfrac{-4\left( {x - \alpha }\right) \left( {x - \beta }\right)  + 3}{2{\left( {x}^{2} + 1\right) }^{2}},$
        因为 $\left( {x - \alpha }\right) \left( {x - \beta }\right)  \leq  0$ ,所以 ${f}^{\prime }\left( x\right)  > 0$ ,
        所以 $f\left( x\right)$ 在 $\left\lbrack  {\alpha ,\beta }\right\rbrack$ 上是单调递增函数,
        所以MAX $f\left( x\right)  = f\left( \beta \right) ,\operatorname{MIN}f\left( x\right)  = f\left( \alpha \right)$ ,
        于是 $g\left( t\right)  = f\left( \beta \right)  - f\left( \alpha \right)  = \dfrac{{2\beta } - t}{{\beta }^{2} + 1} - \dfrac{{2\alpha } - t}{{\alpha }^{2} + 1} = \dfrac{\left( {\beta  - \alpha }\right) \left\lbrack  {t\left( {\alpha  + \beta }\right)  - {2\alpha \beta } + 2}\right\rbrack  }{{\alpha }^{2}{\beta }^{2} + {\alpha }^{2} + {\beta }^{2} + 1} = \dfrac{\left( {\beta  - \alpha }\right) \left\lbrack  {t\left( {\alpha  + \beta }\right)  - {2\alpha \beta } + 2}\right\rbrack  }{{\left( \alpha \beta \right) }^{2} + {\left( \alpha  + \beta \right) }^{2} - {2\alpha \beta } + 1},$
        因为 $\alpha  + \beta  = t,{\alpha \beta } =  - \dfrac{1}{4},\beta  > \alpha$ ,
        所以 $\beta  - \alpha  = \sqrt{{\left( \alpha  + \beta \right) }^{2} - {4\alpha \beta }} = \sqrt{{t}^{2} + 1}$ ,
        将 $\alpha  + \beta  = t,{\alpha \beta } =  - \dfrac{1}{4},\beta  - \alpha  = \sqrt{{t}^{2} + 1}$ 代入 $g\left( t\right)  = \dfrac{\left( {\beta  - \alpha }\right) \left\lbrack  {t\left( {\alpha  + \beta }\right)  - {2\alpha \beta } + 2}\right\rbrack  }{{\left( \alpha \beta \right) }^{2} + {\left( \alpha  + \beta \right) }^{2} - {2\alpha \beta } + 1}$ ,
        得 $g\left( t\right)  = \dfrac{\sqrt{{t}^{2} + 1}\left( {{t}^{2} + \dfrac{5}{2}}\right) }{{t}^{2} + \dfrac{25}{16}} = \dfrac{8\sqrt{{t}^{2} + 1}\left( {2{t}^{2} + 5}\right) }{{16}{t}^{2} + {25}}$ .
    \end{jieda}

    \item 已知函数 $f\left( x\right)  = \ln x - \dfrac{1}{2}{x}^{2}$ . 
    \begin{enumerate}
        \item 求函数 $f\left( x\right)$ 在 $\left\lbrack  {\dfrac{1}{2},4}\right\rbrack$ 上的最大值和最小值
        \item 若不等式 $f\left( x\right)  > \left( {2 - a}\right) {x}^{2}$ 有解,求实数 $a$ 的取值范围
    \end{enumerate}
    \begin{jieda}
        \begin{enumerate}
            \item 因为函数 $f\left( x\right)  = \ln x - \dfrac{1}{2}{x}^{2}$，所以 ${f}^{\prime }\left( x\right)  = \dfrac{1}{x} - x = \dfrac{1 - {x}^{2}}{x}\left( {x > 0}\right)$ .令 ${f}^{\prime }\left( x\right)  = 0$ ,则 $x = 1$ 或 $x =  - 1$ (舍去).
            
            当 $\dfrac{1}{2} < x < 1$ 时, ${f}^{\prime }\left( x\right)  > 0$ ,当 $1 < x < 4$ 时, ${f}^{\prime }\left( x\right)  < 0$
            
            所以函数 $f\left( x\right)$ 在 $\left( {\dfrac{1}{2},1}\right)$ 上单调递增,在$(1,4)$上单调递减，所以当 $x = 1$ 时, $f\left( x\right)$ 取得最大值. 最大值为 $f\left( 1\right)  =  - \dfrac{1}{2}$，又 $f\left( \dfrac{1}{2}\right)  = \ln \dfrac{1}{2} - \dfrac{1}{8} =  - \ln 2 - \dfrac{1}{8},f\left( 4\right)  = 2\ln 2 - 8$
            
            $f\left( \dfrac{1}{2}\right)  - f\left( 4\right)  =  - 3\ln 2 - \dfrac{1}{8} + 8 >  - 3 - \dfrac{1}{8} + 8 > 0.$
            
            所以 $f\left( \dfrac{1}{2}\right)  > f\left( 4\right)$ . 所以当 $x = 4$ 时, $f\left( x\right)$ 取得最小值. 最小值为 $f\left( 4\right)  = 2\ln 2 - 8$
            
            故 $f\left( x\right)$ 在 $\left\lbrack  {\dfrac{1}{2},4}\right\rbrack$ 上的最大值为 $- \dfrac{1}{2}$ . 最小值为 $2\ln 2 - 8$.
            
            \item 易知 $f\left( x\right)$ 的定义域为 $\left( {0, + \infty }\right)$ .故不等式 $f\left( x\right)  > \left( {2 - a}\right) {x}^{2}$ 可化为 $2 - a < \dfrac{\ln x}{{x}^{2}} - \dfrac{1}{2}$.
            
            记 $g\left( x\right)  = \dfrac{\ln x}{{x}^{2}} - \dfrac{1}{2}$ . 则原不等式有解可转化为 $2 - a < g{\left( x\right) }_{\max }$，易得 ${g}^{\prime }\left( x\right)  = \dfrac{1 - 2\ln x}{{x}^{3}},x \in  \left( {0,\sqrt{{e}}}\right)$ 时, ${g}^{\prime }\left( x\right)  > 0,x \in  \left( {\sqrt{{e}}, + \infty }\right)$ 时, ${g}^{\prime }\left( x\right)  < 0$，所以函数 $g\left( x\right)$ 在 $\left( {0,\sqrt{{e}}}\right)$ 上单调递增,在 $\left( {\sqrt{{e}}, + \infty }\right)$ 上单调递减，故 $g{\left( x\right) }_{\max } = g\left( \sqrt{{e}}\right)  = \dfrac{1}{2{e}} - \dfrac{1}{2}$ ,所以 $2 - a < \dfrac{1}{2{e}} - \dfrac{1}{2}$
            
            解得 $a > \dfrac{5}{2} - \dfrac{1}{2{e}}$.
        \end{enumerate}
    \end{jieda}

    \item 已知函数 $f\left( x\right)  = \dfrac{a\ln x}{x} + x - 1$ .
    \begin{enumerate}
        \item 当 $a = 1$ 时，求函数 $g\left( x\right)  = f\left( x\right)  - x$ 极值
        \item 若对任意 $x \in  \lbrack 1, + \infty ),f\left( x\right)  \geq  a + 1$ 恒成立,求实数 $a$ 的取值范围
    \end{enumerate}
    \begin{jieda}
        \begin{enumerate}
            \item 函数 $f\left( x\right)  = \dfrac{a\ln x}{x} + x - 1$ 的定义域为 $\left( {0, + \infty }\right)$ . 当 $a = 1$ 时, $g\left( x\right)  = f\left( x\right)  - x = \frac{\ln x}{x} - 1$.
            求导得 ${g}^{\prime }\left( x\right)  = \dfrac{1 - \ln x}{{x}^{2}}$ . 由 ${g}^{\prime }\left( x\right)  = 0$ . 得 $x = {e}$ . 由 ${g}^{\prime }\left( x\right)  > 0$ . 得 $0 < x < {e}$ . 由 ${g}^{\prime }\left( x\right)  < 0$ ,得 $x > {e}$ .
            因此 $g\left( x\right)$ 在$(0,e)$上单调递增. 在 $\left( {{e}, + \infty }\right)$ 上单调递减.
            所以 $g\left( x\right)$ 在 $x = {e}$ 处取得极大值 $g\left( {e}\right)  = \dfrac{1}{{e}} - 1$ ,无极小值.
            \item 函数 $f\left( x\right)  = \dfrac{a\ln x}{x} + x - 1$ , $f\left( x\right)  \geq  a + 1 \Longleftrightarrow  a\left( {\ln x - x}\right)  \geq  {2x} - {x}^{2}$ , $x \in  \lbrack 1, + \infty )$ .
            设 $m\left( x\right)  = \ln x - x.x \in  \lbrack 1, + \infty )$ . 求导得 ${m}^{\prime }\left( x\right)  = \dfrac{1 - x}{x} \leq  0$ . 函数 $m\left( x\right)$ 在 $\lbrack 1, + \infty )$ 上单调递减.
            则 $m\left( x\right)  \leq  m\left( 1\right)  =  - 1 < 0$ . 即 $\ln x < x$ . 因此 $a \leq  \dfrac{{x}^{2} - {2x}}{x - \ln x}$ .
            令 $\varphi \left( x\right)  = \dfrac{{x}^{2} - {2x}}{x - \ln x},x \in  \lbrack 1, + \infty )$ . 求导得 ${\varphi }^{\prime }\left( x\right)  = \dfrac{\left( {x - 1}\right) \left( {x + 2 - 2\ln x}\right) }{{\left( x - \ln x\right) }^{2}}$ .
            令 $h\left( x\right)  = x + 2 - 2\ln x,x \in  \lbrack 1, + \infty )$ ,求导得 ${h}^{\prime }\left( x\right)  = 1 - \dfrac{2}{x}$ . 当 $1 \leq  x < 2$ 时, ${h}^{\prime }\left( x\right)  < 0$ ,
            当 $x > 2$ 时, ${h}^{\prime }\left( x\right)  > 0$ ,即 $h\left( x\right)$ 在(1,2)上单调递减,在 $\left( {2, + \infty }\right)$ 上单调递增.
            则 $h{\left( x\right) }_{\min } = h\left( 2\right)  = 4 - 2\ln 2 > 0$ ,即 ${\varphi }^{\prime }\left( x\right)  \geq  0$ ,因此函数 $\varphi \left( x\right)$ 在 $\lbrack 1, + \infty )$ 上是增函数.
            $\varphi {\left( x\right) }_{\min } = \varphi \left( 1\right)  =  - 1$
            所以 $a \leq   - 1$ ,即实数 $a$ 的取值范围为 $( - \infty , - 1\rbrack$ .
        \end{enumerate}
    \end{jieda}
\end{enumerate}